\documentclass[11pt]{article}

\usepackage[T1]{fontenc}
\usepackage{amsmath,amssymb,amsthm}
\usepackage{geometry}
\usepackage{hyperref}

\geometry{margin=1in}
\hypersetup{colorlinks=true,linkcolor=blue,citecolor=blue,urlcolor=blue}

\theoremstyle{plain}
\newtheorem{theorem}{Theorem}[section]
\newtheorem{proposition}[theorem]{Proposition}
\newtheorem{lemma}[theorem]{Lemma}
\newtheorem{corollary}[theorem]{Corollary}

\theoremstyle{definition}
\newtheorem{definition}[theorem]{Definition}

\theoremstyle{remark}
\newtheorem{remark}[theorem]{Remark}

\newcommand{\Vir}{\operatorname{Vir}}
\newcommand{\Tr}{\operatorname{Tr}}
\newcommand{\rad}{\operatorname{rad}}
\newcommand{\ord}{\operatorname{ord}}
\newcommand{\wt}{\operatorname{wt}}
\newcommand{\Res}{\operatorname{Res}}

\title{Generalized Rocha--Caridi Identities and Exact Confluence\\
for Irreducible Minimal-Model Torus One-Point Traces}
\author{Leslie P. Polzer\\
\small Independent Researcher\\
\small \texttt{polzer@fastmail.com}}
\date{November 20, 2026}

\begin{document}
\maketitle

\begin{abstract}
Let \(b_\star^2=-p'/p\) (\(p,p'\geq2\) coprime) be a Virasoro minimal-model
charge, let \(\lambda=h_{r,s}(b_\star)\) be an interior Kac weight, and let
\(L(c_\star,\lambda)=M(c_\star,\lambda)/J(c_\star,\lambda)\) be the
irreducible quotient of the Verma module.  We prove three exact algebraic
theorems for the irreducible torus one-point trace
\(\mathcal{T}^{\mathrm{irr}}_{c_\star,\lambda,d}(q)
=\Tr_{L(c_\star,\lambda)}(q^{L_0-\lambda}V_d(1))\).
First, the three-point form \(\rho(c_\star,\lambda,d)\) descends to
\(L(c_\star,\lambda)\times L(c_\star,\lambda)\) if and only if \(d\) belongs to
the minimal-model fusion set \(\mathcal{F}_{p,p'}(r,s)\), whereas the
Verma-module torus one-point block coefficients \(F_N(c_\star,\Delta,d)\) are
regular (pole order zero) at \(\Delta=\lambda\) for every \(N\geq0\) if and
only if \(d\) belongs to the reflected set
\(\mathcal{F}^\pm_{p,p'}(r,s)\supseteq\mathcal{F}_{p,p'}(r,s)\); for every
\(d\in\mathcal{F}_{p,p'}(r,s)\), all \(j_N(\lambda)\) Jantzen poles cancel
identically through a Jantzen vertex filtration on \(M(c_\star,\lambda+x)[[x]]\).
Second, for generic \(b\) and any Kac label \((R,S)\) with fusion insertion
\(d=h_{2k+1,2m+1}(b)\) (\(0\leq k<R\), \(0\leq m<S\)), we prove via a
two-weight fusion product rule that the first-layer Jantzen vertex ratio
\(c_{R,S}(b^2;2k+1,2m+1)\) equals an explicit shifted rectangle product that
factorizes as a ratio of Gamma potentials \(\Psi_{2k+1,2m+1}\).
Third, specializing the generic degenerate quotient \(V_{R_1,S_1}(b_\star+t)\)
yields an exact length-two \(t\)-Jantzen identity that telescopes along the
braided BGG orbit \(\mathcal{O}_{p,p'}(r,s)\) to the generalized
Rocha--Caridi identity
\(\mathcal{T}^{\mathrm{irr}}_{c_\star,\lambda,d}(q)
=\mathcal{F}^{\mathrm{Verma}}_{c_\star,\lambda,d}(q)
+\sum_{j\geq1}(-1)^j\sum_{\pm}\Omega_{u,v}(\lambda\to\mu^\pm_j)\,
q^{\ell^\pm_j}\mathcal{F}^{\mathrm{Verma}}_{c_\star,\mu^\pm_j,d}(q)\),
with explicit nonzero rational weights
\(\Omega_{u,v}(\lambda\to\mu^\pm_j)\in\mathbb{Q}^\times\).
Combined with fixed-\(b_\star\) algebraic confluence for Verma torus blocks,
this computes every coefficient of
\(\mathcal{T}^{\mathrm{irr}}_{c_\star,\lambda,d}(q)\) exactly from the
generic Poghossian--Hadasz--Jaskolski--Suchanek recursion with no numerical
\(\varepsilon\)-limit; for single-index insertions \(d=h_{1,2m+1}(b_\star)\)
or \(d=h_{2k+1,1}(b_\star)\), the truncated recursion is completely
collision-free and pole-free at \(b=b_\star\) to all orders.
\end{abstract}

\section{Introduction}\label{sec:intro}

In a Virasoro minimal model \(\mathcal{M}(p,p')\) with coprime integers
\(p,p'\geq2\) and central charge \(c_\star=1-6(p'-p)^2/(pp')\), the physical
highest-weight representations are the irreducible quotients
\(L(c_\star,\lambda)=M(c_\star,\lambda)/J(c_\star,\lambda)\) of Verma modules
at interior Kac weights \(\lambda=h_{r,s}(b_\star)\) (\(0<r<p\), \(0<s<p'\),
\(b_\star^2=-p'/p\)).  When the external operator is the identity (\(d=0\)),
the torus trace is the minimal-model character, expressed by Rocha--Caridi
\cite{RochaCaridi1985} as an alternating sum of Verma characters over the
braided Feigin--Fuchs embedding diagram of \(M(c_\star,\lambda)\).  For a
general primary insertion \(V_d(1)\) of weight \(d\neq0\), the irreducible
torus one-point trace
\begin{equation}\label{eq:irr-trace-def}
 \mathcal{T}^{\mathrm{irr}}_{c_\star,\lambda,d}(q)
 =\Tr_{L(c_\star,\lambda)}\bigl(q^{L_0-\lambda}V_d(1)\bigr)
 =\sum_{N=0}^\infty T^{\mathrm{irr}}_N(c_\star,\lambda,d)\,q^N
\end{equation}
was previously accessed either through Coulomb-gas screening contour integrals
in Felder's Fock-space BRST complex \cite{Felder1989,FelderSilvotti1989} or
through direct level-by-level inversion of the quotient Shapovalov matrix.
From the recursive side, Poghossian \cite{Poghossian2009} and
Hadasz--Jaskolski--Suchanek \cite{HJS2010} derived an all-order residue
recursion for the generic-\(b\) Verma torus one-point block
\(\mathcal{F}^{\mathrm{Verma}}_{c,\Delta,d}(q)
=\sum_{N\geq0}F_N(c,\Delta,d)q^N\).  At rational \(b_\star^2\), individual
summands develop non-physical poles; Javerzat, Santachiara, and Foda
\cite[Section~5]{JSF2018} recovered the \(d=0\) minimal-model characters from
the generic torus recursion via a fine three-parameter limit (verifying the
resulting rational chain sum numerically up to order \(300\)), while Ribault
\cite[Section~3]{Ribault2019} and \cite[Section~4.3.3]{Ribault2026} framed a
fixed-rational-\(b^2\) recursion as an open problem.

In the companion work \cite{CompanionTorus}, exact fixed-\(b_\star\) algebraic
confluence and sharp Jantzen pole orders were proved for the \emph{Verma-module}
torus one-point block: for indeterminate \(d\), \(F_N(c_\star,\Delta,d)\) has a
pole of exact order \(j_N(\lambda)\) at \(\Delta=\lambda\), growing without
bound as \(N\to\infty\).  The present paper solves the corresponding problems
for the \emph{irreducible} minimal-model torus one-point traces
\eqref{eq:irr-trace-def}.

A direct small-case computation reveals immediately why naive specialization
fails.  For the Ising model \((p,p')=(3,4)\) at \(\lambda=h_{1,2}=1/16\) and
\(d=h_{1,3}=1/2\), the Verma block \(F_N(1/2,\Delta,1/2)\) is regular at
\(\Delta=1/16\) at every level, yet its value at \(\Delta=1/16\) does not equal
the irreducible trace: at \(N=2\) one has \(F_2(1/2,1/16,1/2)=8/7\) whereas
\(T^{\mathrm{irr}}_2(1/2,1/16,1/2)=-1\), a discrepancy of \(15/7\).  Moving
along a fixed-\(c_\star\) ray \((\Delta,d)=(\lambda+x,d_\star+\alpha x)\)
removes the discrepancy at \(N=2,3\) for \(\alpha=5\), but fails at \(N=4\)
(which requires \(\alpha=17/19\)), and three-parameter rays vary continuously
with their slopes.  The discrepancy comes from the first Jantzen layer
\(J^1/J^2\): because \(d=1/2\) kills the vertex form \(\rho\) on \(J^1\),
the \(O(x)\) vanishing of \(\rho_x\) cancels the \(O(x)\) vanishing of the
Shapovalov form \(B_x\) on \(J^1\), leaving a finite \(O(x^0)\) contribution
\(\sigma^\pm_1/\gamma^\pm_1\) that survives the limit \(\Delta\to\lambda\).

Tracking this mechanism across the full braided embedding diagram yields three
main theorems.

\begin{theorem}[Fusion-descent and exact Verma regularity]\label{thm:descent-regularity}
Let \(p,p'\geq2\) be coprime, \(b_\star^2=-p'/p\), and
\(\lambda=h_{r,s}(b_\star)\) with \(0<r<p\) and \(0<s<p'\).  Define the
minimal-model fusion set and its reflected enlargement
\begin{align}
 \mathcal{F}_{p,p'}(r,s)
 &:=\bigl\{\,h_{2k+1,2m+1}(b_\star)
   : 0\leq k<\min(r,p-r),\;\; 0\leq m<\min(s,p'-s)\,\bigr\},
   \label{eq:fusion-set}\\
 \mathcal{F}^\pm_{p,p'}(r,s)
 &:=\bigl\{\,h_{2k+1,\pm(2m+1)}(b_\star)
   : 0\leq k<\min(r,p-r),\;\; 0\leq m<\min(s,p'-s)\,\bigr\}.
   \label{eq:fusion-set-pm}
\end{align}
Then the map \((k,m)\mapsto h_{2k+1,2m+1}(b_\star)\) is injective on the index
box of \eqref{eq:fusion-set}, and for every \(d\in\mathbb{C}\):
\begin{enumerate}
 \item[(i)] the Verma three-point form
   \(\rho(c_\star,\lambda,d):M(c_\star,\lambda)\times M(c_\star,\lambda)\to\mathbb{C}\)
   descends to a well-defined bilinear form on
   \(L(c_\star,\lambda)\times L(c_\star,\lambda)\) (so \(V_d(1)\) descends to a
   primary operator on \(L(c_\star,\lambda)\) and
   \(\mathcal{T}^{\mathrm{irr}}_{c_\star,\lambda,d}(q)\) is well-defined) if
   and only if \(d\in\mathcal{F}_{p,p'}(r,s)\);
 \item[(ii)] the Verma torus one-point block coefficient
   \(F_N(c_\star,\Delta,d)\) is regular (pole order zero) at \(\Delta=\lambda\)
   for every \(N\geq0\) if and only if \(d\in\mathcal{F}^\pm_{p,p'}(r,s)\).
\end{enumerate}
When \(d=h_{2k+1,\epsilon(2m+1)}(b_\star)\in\mathcal{F}^\pm_{p,p'}(r,s)\)
(\(\epsilon\in\{+,-\}\)), \(F_N(c_\star,\Delta,d)\) is regular at every weight
\(\mu\) in the braided BGG orbit \(\mathcal{O}_{p,p'}(r,s)\) of \(\lambda\),
and \(F_N(c(b),\Delta,h_{2k+1,\epsilon(2m+1)}(b))\) is jointly regular in
\((b,\Delta)\) at \((b_\star,\mu)\) for every \(N\geq0\).
\end{theorem}

To isolate \(\mathcal{T}^{\mathrm{irr}}\) from \(\mathcal{F}^{\mathrm{Verma}}\),
we evaluate the first-layer Jantzen vertex ratio at generic \(b\).

\begin{theorem}[Generic first-layer Jantzen vertex ratio]\label{thm:first-layer-ratio}
Let \(b\in\mathbb{C}^\times\) be generic, put \(t=b^2\), fix positive integers
\(R,S\geq1\), and let \(0\leq k\leq R-1\) and \(0\leq m\leq S-1\) with
\(d=h_{2k+1,2m+1}(b)\).  Let \(\chi_{R,S}(b)\in M(c(b),h_{R,S}(b))_{RS}\) be
the nonzero singular vector (with constant PBW coefficients lifted to
\(M(c(b),h_{R,S}(b)+x)\)).  Then
\begin{equation}\label{eq:box-formula}
 c_{R,S}(t;2k+1,2m+1)
 :=\frac{\left.\partial_\Delta\right|_{\Delta=h_{R,S}(b)}
         \rho_{\Delta,d}(\chi_{R,S},\chi_{R,S})}
        {\left.\partial_\Delta\right|_{\Delta=h_{R,S}(b)}
         B_\Delta(\chi_{R,S},\chi_{R,S})}
 =\frac{\displaystyle\prod_{\substack{1-R\leq a\leq R\\1-S\leq b'\leq S\\(a,b')\neq(-k,-m)}}
        \bigl((a+k)t+(b'+m)\bigr)}
       {\displaystyle\prod_{\substack{1-R\leq a\leq R\\1-S\leq b'\leq S\\(a,b')\neq(0,0)}}
        (at+b')}.
\end{equation}
Moreover, setting \(P_{R,S}=(Rb+Sb^{-1})/2\) and
\(P_{R,-S}=(Rb-Sb^{-1})/2\), this rational function factorizes as
\(c_{R,S}(t;2k+1,2m+1)=\Psi_{2k+1,2m+1}(P_{R,-S};b)/\Psi_{2k+1,2m+1}(P_{R,S};b)\),
where
\begin{equation}\label{eq:gamma-potential}
 \begin{aligned}
  \Psi_{2k+1,2m+1}(P;b)
  :=\prod_{j=1}^m&\frac{1}{\Gamma\bigl(1+\tfrac{2P}{b}+\tfrac{j}{b^2}\bigr)
                           \Gamma\bigl(1-\tfrac{2P}{b}+\tfrac{j}{b^2}\bigr)}\\
   \times\prod_{i=1}^k&\frac{1}{\Gamma(1+2Pb+ib^2+m)\,\Gamma(1-2Pb+ib^2+m)}.
 \end{aligned}
\end{equation}
\end{theorem}

The proof of Theorem~\ref{thm:first-layer-ratio} is completely elementary: it
introduces the symmetric two-weight three-point form
\(\rho_{\Delta_1,\Delta_2;d}\) and differentiates its factorization along
\(\Delta_2=h_{R,S}(b)\), dividing by the \(d=0\) specialization
\(\rho_{\Delta,\Delta;0}\equiv B_\Delta\).  No free-field integrals or AGT
machinery are used.

Our third theorem expresses every irreducible minimal-model torus one-point
trace as an alternating BGG sum of regular Verma torus one-point blocks.  Let
\(\{\mu^\pm_j\}_{j\geq1}\) be the BGG orbit weights of
\(\lambda=h_{r,s}(b_\star)\) at levels \(\ell^\pm_j=\mu^\pm_j-\lambda\)
(Section~\ref{sec:setup}).  For each primitive BGG arrow
\(\mu\to\nu\) of Kac label \((R,S)\), every factor in the reduced box ratio
\(c_{R,S}(-p'/p;2k+1,2m+1)\) is nonzero, and the product of box ratios along
any BGG path from \(\lambda\) to \(\mu^\pm_j\) is path-independent, defining a
canonical rational weight
\(\Omega_{2k+1,2m+1}(\lambda\to\mu^\pm_j)\in\mathbb{Q}^\times\)
(Proposition~\ref{prop:bgg-cocycle}).

\begin{theorem}[Generalized Rocha--Caridi identity and exact confluence]\label{thm:gen-rocha-caridi}
Let \(b_\star^2=-p'/p\) (\(p,p'\geq2\) coprime), let
\(\lambda=h_{r,s}(b_\star)\) be an interior Kac weight, and let
\(d=h_{u,v}(b_\star)\in\mathcal{F}_{p,p'}(r,s)\) with
\((u,v)=(2k+1,2m+1)\).  Then as formal power series in \(q\) (with finite sums
at every coefficient \(q^N\)),
\begin{equation}\label{eq:gen-rocha-caridi}
 \mathcal{T}^{\mathrm{irr}}_{c_\star,\lambda,d}(q)
 =\mathcal{F}^{\mathrm{Verma}}_{c_\star,\lambda,d}(q)
  +\sum_{j=1}^\infty(-1)^j\Bigl(
    \Omega_{u,v}(\lambda\to\mu^+_j)\,q^{\ell^+_j}\,
    \mathcal{F}^{\mathrm{Verma}}_{c_\star,\mu^+_j,d}(q)
   +\Omega_{u,v}(\lambda\to\mu^-_j)\,q^{\ell^-_j}\,
    \mathcal{F}^{\mathrm{Verma}}_{c_\star,\mu^-_j,d}(q)\Bigr).
\end{equation}
Consequently:
\begin{enumerate}
 \item[(a)] At \(d=0\) (\(k=m=0\)), \(\Omega_{1,1}\equiv1\) and
   \(\mathcal{F}^{\mathrm{Verma}}_{c_\star,\mu,0}(q)=\phi(q)^{-1}\),
   recovering the Rocha--Caridi character formula \cite{RochaCaridi1985}.
 \item[(b)] For every \(N\geq0\), combining \eqref{eq:gen-rocha-caridi} with
   fixed-\(b_\star\) Verma confluence \cite[Theorem~4.3]{CompanionTorus}
   computes \(T^{\mathrm{irr}}_N(c_\star,\lambda,d)\in\mathbb{Q}\) exactly from
   the generic Poghossian--HJS recursion by a finite algebraic procedure with
   no numerical \(\varepsilon\)-limit.
 \item[(c)] For every layer \(j\geq1\), the \(j\)-th Jantzen vertex ratio on
   \(M(c_\star,\lambda)\) satisfies
   \(\sigma^\pm_j(\lambda)/\gamma^\pm_j(\lambda)
   =\Omega_{u,v}(\lambda\to\mu^\pm_j)\).
\end{enumerate}
\end{theorem}

Remarkably, for single-index fusion insertions (which include every nontrivial
primary in Lee--Yang \((2,5)\), Ising \((3,4)\), \((2,7)\), \((3,5)\), and
both \(d=3/5\) and \(d=3/2\) in Tricritical Ising \((4,5)\)), even the
off-orbit collisions vanish completely.

\begin{corollary}[Collision-free single-index recursion]\label{cor:collision-free}
In Theorem~\ref{thm:gen-rocha-caridi}, suppose \(d=h_{1,2m+1}(b_\star)\) with
\(0\leq m<\min(s,p'-s)\) (or dually \(d=h_{2k+1,1}(b_\star)\) with
\(0\leq k<\min(r,p-r)\)).  Then in the truncated Poghossian--HJS recursion for
\(H_N(c(b),\Delta,d(b))\) restricted to Kac labels \((R,S)\) with \(S\leq m\)
(resp.\ \(R\leq k\)), every residue \(R_{R,S}(b,d(b))\) is regular at
\(b=b_\star\), no two active poles \(h_{R_1,S_1}(b_\star)\) and
\(h_{R_2,S_2}(b_\star)\) collide, and no shifted weight
\(h_{R_1,-S_1}(b_\star)\) hits any active pole \(h_{R_2,S_2}(b_\star)\).  Thus
every \(\mathcal{F}^{\mathrm{Verma}}_{c_\star,\mu^\pm_j,d}(q)\) in
\eqref{eq:gen-rocha-caridi} is evaluated directly at \(b=b_\star\) by the
unmodified truncated recursion, with no formal \(t\)-expansion at all.
\end{corollary}

The proof of Theorem~\ref{thm:gen-rocha-caridi} bypasses all direct
higher-layer Jantzen computations.  For each weight \(\mu\) in the BGG orbit
with primitive singular vectors of labels \((R_1,S_1)\) and \((R_2,S_2)\)
leading to \(\nu_1,\nu_2\), we deform \(b=b_\star+t\) along the generic
degenerate quotient
\[
 V_{R_1,S_1}(b):=M\bigl(c(b),h_{R_1,S_1}(b)\bigr)\big/
                 M\bigl(c(b),h_{R_1,-S_1}(b)\bigr).
\]
At \(t=0\), its specialization \(V_{R_1,S_1}(b_\star)=M(c_\star,\mu)/M(c_\star,\nu_1)\)
has exact length two with composition factors \(L(c_\star,\mu)\) and
\(L(c_\star,\nu_2)\), and its \(t\)-Jantzen filtration terminates at depth one
(\(J_t^2=0\)).  Taking \(t\to0\) yields a two-term identity relating
\(\mathcal{T}^{\mathrm{irr}}(\mu)\) and \(\mathcal{T}^{\mathrm{irr}}(\nu_2)\)
to \(\mathcal{F}^{\mathrm{Verma}}(\mu)\) and
\(\mathcal{F}^{\mathrm{Verma}}(\nu_1)\), which telescopes along the BGG chain
to \eqref{eq:gen-rocha-caridi}.

\section{Setup, Two-Weight Ward Identities, and Braided Orbits}\label{sec:setup}

Let \(\Vir=\bigoplus_{n\in\mathbb{Z}}\mathbb{C}L_n\oplus\mathbb{C}C\) be the
Virasoro algebra with
\([L_m,L_n]=(m-n)L_{m+n}+\frac{m^3-m}{12}\delta_{m+n,0}C\).  For
\((c,\Delta)\in\mathbb{C}^2\), the Verma module
\(M(c,\Delta)=\bigoplus_{N\geq0}M(c,\Delta)_N\) has highest-weight vector
\(v_\Delta\) (\(L_0v_\Delta=\Delta v_\Delta\), \(L_mv_\Delta=0\) for
\(m>0\)) and PBW basis \(\{L_{-Y}v_\Delta:|Y|=N\}\) at level \(N\), where
\(Y=(y_1\geq\cdots\geq y_\ell\geq1)\) ranges over partitions of \(N\) and
\(L_{-Y}=L_{-y_1}\cdots L_{-y_\ell}\).  The contravariant Shapovalov form
\(B_\Delta=B_{c,\Delta}:M(c,\Delta)\times M(c,\Delta)\to\mathbb{C}\) is the
unique symmetric bilinear form with \(B_\Delta(v_\Delta,v_\Delta)=1\) and
\(B_\Delta(L_{-n}u,v)=B_\Delta(u,L_nv)\).  Its level-\(N\) Gram matrix is
\(G_N(c,\Delta)_{Y,Y'}=B_\Delta(L_{-Y}v_\Delta,L_{-Y'}v_\Delta)\).  Its
radical \(J(c,\Delta)=\rad B_\Delta\) is the unique maximal proper submodule
of \(M(c,\Delta)\), and \(L(c,\Delta)=M(c,\Delta)/J(c,\Delta)\) is the
irreducible quotient, on which the induced form \(\bar B\) is nondegenerate.

For \(b\in\mathbb{C}^\times\) we set \(t=b^2\), \(Q=b+b^{-1}\),
\(c(b)=1+6Q^2=13+6t+6t^{-1}\), and for any integers \(r,s\in\mathbb{Z}\),
\begin{equation}\label{eq:kac-weight}
 h_{r,s}(b)
 :=\frac{Q^2-(rb+sb^{-1})^2}{4}
  =\frac{(t+1)^2-(rt+s)^2}{4t},
\end{equation}
so that \(h_{r,-s}(b)=h_{r,s}(b)+rs\).  When \(b_\star^2=-p'/p\) with
coprime integers \(p,p'\geq2\), we have \(c_\star=c(b_\star)=1-6(p'-p)^2/(pp')\)
and
\begin{equation}\label{eq:mm-weight}
 h_{r,s}(b_\star)
 =\frac{(p'r-ps)^2-(p'-p)^2}{4pp'}.
\end{equation}

\subsection{Two-weight three-point pairings and fusion polynomials}

To prove Theorem~\ref{thm:first-layer-ratio}, we work with the bilinear form
between two Verma modules of possibly distinct highest weights.

\begin{lemma}[Two-weight primary Ward identity]\label{lem:two-weight-ward}
Fix \(c,\Delta_1,\Delta_2,d\in\mathbb{C}\).  There is a unique bilinear form
\(\rho_{\Delta_1,\Delta_2;d}:M(c,\Delta_1)\times M(c,\Delta_2)\to\mathbb{C}\)
such that \(\rho_{\Delta_1,\Delta_2;d}(v_{\Delta_1},v_{\Delta_2})=1\) and,
for every \(m\in\mathbb{Z}\) and homogeneous vectors
\(u\in M(c,\Delta_1)_{N_1}\), \(v\in M(c,\Delta_2)_{N_2}\),
\begin{equation}\label{eq:two-weight-ward}
 \rho_{\Delta_1,\Delta_2;d}(L_{-m}u,v)
 -\rho_{\Delta_1,\Delta_2;d}(u,L_mv)
 =\bigl((\Delta_1+N_1)-(\Delta_2+N_2)+md\bigr)\,
  \rho_{\Delta_1,\Delta_2;d}(u,v).
\end{equation}
Its matrix entries in the PBW bases belong to
\(\mathbb{Q}[c,\Delta_1,\Delta_2,d]\), it satisfies the symmetry
\begin{equation}\label{eq:two-weight-sym}
 \rho_{\Delta_1,\Delta_2;d}(u,v)=\rho_{\Delta_2,\Delta_1;d}(v,u),
\end{equation}
and on the diagonal \(\Delta_1=\Delta_2=\Delta\) at \(d=0\) it coincides with
the Shapovalov form:
\begin{equation}\label{eq:d0-shapovalov}
 \rho_{\Delta,\Delta;0}(u,v)=B_{c,\Delta}(u,v)
 \qquad\text{for all }u,v\in M(c,\Delta).
\end{equation}
\end{lemma}

\begin{proof}
The relation \eqref{eq:two-weight-ward} is the \(z=1\) specialization of the
primary Ward identity
\([L_m,V_d(z)]=z^m(z\partial_z+(m+1)d)V_d(z)=z^m([L_0,V_d(z)]+mdV_d(z))\)
acting on \(\langle u|V_d(z)|v\rangle=z^{\wt(u)-\wt(v)-d}\rho(u,v)\).  Induction
on \(N_1+N_2\) computes \(\rho_{\Delta_1,\Delta_2;d}(L_{-Y}v_{\Delta_1},
L_{-Y'}v_{\Delta_2})\) uniquely by moving positive modes from left to right
(using \(m>0\)) and then moving negative modes from right to left (using
\(m<0\) and \(L_{-m}v_{\Delta_1}=0\) for \(m<0\)), producing polynomials in
\(\mathbb{Q}[c,\Delta_1,\Delta_2,d]\).  Existence of this coefficientwise formal bilinear pairing is standard from the
normalized three-Verma three-point block
\cite[Section~3, p.~4]{HJS2010} (equivalently, the \(z=1\) evaluation of the
unique \(\Vir\)-covariant formal pairing into
\(z^{\Delta_1-\Delta_2-d}\mathbb{C}[z,z^{-1}]\)).  Swapping
\((\Delta_1,u)\leftrightarrow(\Delta_2,v)\) and replacing \(m\) by \(-m\) in
\eqref{eq:two-weight-ward} negates both sides, leaving the defining identity
unchanged; since the value on \((v_{\Delta_1},v_{\Delta_2})\) is \(1\),
uniqueness gives \eqref{eq:two-weight-sym}.  Finally, when
\(\Delta_1=\Delta_2=\Delta\) and \(d=0\), the Shapovalov form \(B_{c,\Delta}\)
satisfies \(B_{c,\Delta}(v_\Delta,v_\Delta)=1\), \(B_{c,\Delta}(u,v)=0\) for
\(N_1\neq N_2\), and \(B_{c,\Delta}(L_{-m}u,v)=B_{c,\Delta}(u,L_mv)\), so
\eqref{eq:two-weight-ward} holds identically for \(B_{c,\Delta}\) (both sides
vanish identically when \(N_1=N_2\) as well as when \(N_1\neq N_2\));
uniqueness gives \eqref{eq:d0-shapovalov}.
\end{proof}

We abbreviate \(\rho_{\Delta,d}:=\rho_{\Delta,\Delta;d}\), and denote its
level-\(N\) matrix by \(\rho_N(c,\Delta,d)\).  The Verma torus one-point block
coefficients are \(F_N(c,\Delta,d):=\Tr(G_N(c,\Delta)^{-1}\rho_N(c,\Delta,d))\)
(noting \(\rho_N^{\mathsf{T}}=\rho_N\) by \eqref{eq:two-weight-sym}), with
generating series \(\mathcal{F}^{\mathrm{Verma}}_{c,\Delta,d}(q)
=\sum_{N\geq0}F_N(c,\Delta,d)q^N\) and eta-reduced block
\(H(c,\Delta,d;q)=\phi(q)\mathcal{F}^{\mathrm{Verma}}_{c,\Delta,d}(q)\), where
\(\phi(q)=\prod_{m=1}^\infty(1-q^m)\).

\begin{proposition}[Singular vectors and fusion polynomials]\label{prop:fusion-poly}
Fix positive integers \(R,S\geq1\).
\begin{enumerate}
 \item[(i)] There is a unique family of singular vectors
   \(\chi_{R,S}(b)\in M(c(b),h_{R,S}(b))_{RS}\) whose coefficient on
   \(L_{-1}^{RS}v\) equals \(1\); its PBW coefficients belong to
   \(\mathbb{Q}[b^2,b^{-2}]\), so \(\chi_{R,S}(b_\star)\) is well-defined and
   nonzero at every \(b_\star\in\mathbb{C}^\times\).
 \item[(ii)] For every \(b\in\mathbb{C}^\times\) and
   \(\Delta_1,\Delta_2\in\mathbb{C}\), writing
   \(\Delta_i=(Q^2-\lambda_i^2)/4\) (\(i=1,2\)),
   \begin{equation}\label{eq:fusion-poly}
    \rho_{h_{R,S}(b),\Delta_1;\Delta_2}\bigl(\chi_{R,S}(b),v_{\Delta_1}\bigr)
    =C_{R,S}(b)\,P_{R,S}(b;\Delta_1,\Delta_2),
   \end{equation}
   where \(C_{R,S}(b)\in\mathbb{C}^\times\) is independent of
   \((\Delta_1,\Delta_2)\) and
   \begin{equation}\label{eq:P-RS-def}
    P_{R,S}(b;\Delta_1,\Delta_2)
    :=\prod_{\substack{p'\in\{1-R,3-R,\ldots,R-1\}\\q'\in\{1-S,3-S,\ldots,S-1\}}}
      \frac{\lambda_1+\lambda_2+p'b+q'b^{-1}}{2}\,
      \frac{\lambda_1-\lambda_2+p'b+q'b^{-1}}{2}.
   \end{equation}
   In particular, at \(\Delta_1=h_{R,S}(b)\) (where \(\lambda_1=Rb+Sb^{-1}\)),
   the roots of \(P_{R,S}(b;h_{R,S}(b),d)=0\) in \(d\) are precisely
   \begin{equation}\label{eq:P-RS-roots}
    d=h_{u,v}(b), \qquad
    u\in\{1,3,\ldots,2R-1\},\quad v\in\{1,3,\ldots,2S-1\}.
   \end{equation}
\end{enumerate}
\end{proposition}

\begin{proof}
Part~(i) is \cite[Theorem~3.1]{Astashkevich1997} (composing the polynomial PBW
coefficients in \((c,h)\) with their Laurent expressions in \(t=b^2\); see also
\cite[Propositions~5.1--5.2 and Corollary~5.2]{IoharaKoga2011} and
\cite[Proposition~6.2 and Lemma~6.6]{CompanionTorus}).  Part~(ii) is
Zamolodchikov's fusion polynomial \cite[Section~5,
equations~(5.9)--(5.12)]{Zamolodchikov2004} (see also \cite[Section~3,
equations~(7)--(9)]{HJS2010}); the root list \eqref{eq:P-RS-roots} follows
directly by setting \(\lambda_1=Rb+Sb^{-1}\) and noting that
\(\lambda_2=\pm(\lambda_1+p'b+q'b^{-1})\) ranges over
\(\pm(ub+vb^{-1})\) with \(u=R+p'\in\{1,3,\ldots,2R-1\}\) and
\(v=S+q'\in\{1,3,\ldots,2S-1\}\).
\end{proof}

\subsection{The braided BGG orbit at minimal-model charges}

Fix coprime \(p,p'\geq2\), \(b_\star^2=-p'/p\), and an interior Kac weight
\(\lambda=h_{r,s}(b_\star)\) (\(0<r<p\), \(0<s<p'\)).  Replacing \((r,s)\) by
\((p-r,p'-s)\) if necessary (which preserves \(h_{r,s}(b_\star)\)), we normalize
so that \(n_1:=p'r+ps<pp'\), and set \(n_0:=p'r-ps\neq0\).  Define the sequence
of integers
\begin{equation}\label{eq:bgg-momenta}
 \begin{aligned}
  n^+_{2j}&:=2jpp'+n_0, &\qquad n^-_{2j}&:=2jpp'-n_0 \qquad (j\geq0,\;n^\pm_0=\pm n_0),\\
  n^+_{2j+1}&:=2jpp'+n_1, &\qquad n^-_{2j+1}&:=2(j+1)pp'-n_1 \qquad (j\geq0),
 \end{aligned}
\end{equation}
and the associated BGG orbit weights and levels
\begin{equation}\label{eq:bgg-weights}
 \mu^\pm_k:=\frac{(n^\pm_k)^2-(p'-p)^2}{4pp'}, \qquad
 \ell^\pm_k:=\mu^\pm_k-\lambda=\frac{(n^\pm_k)^2-n_0^2}{4pp'} \qquad (k\geq1),
\end{equation}
with \(\mu_0:=\lambda\) and \(\mathcal{O}_{p,p'}(r,s):=\{\mu_0\}\cup
\{\mu^\pm_k:k\geq1\}\).  Explicitly,
\(\ell^\pm_{2j}=j(jpp'\pm n_0)\), \(\ell^+_{2j+1}=(jp+r)(jp'+s)\), and
\(\ell^-_{2j+1}=((j+1)p-r)((j+1)p'-s)\).

\begin{lemma}[Structure and arithmetic of primitive BGG arrows]\label{lem:bgg-arrows}
Let \(\mu\in\mathcal{O}_{p,p'}(r,s)\) lie at layer \(k\geq0\) (so
\(\mu=\mu_0\) if \(k=0\), and \(\mu=\mu^\epsilon_k\) with
\(\epsilon\in\{+,-\}\) if \(k\geq1\)).
\begin{enumerate}
 \item[(i)] The Verma module \(M(c_\star,\mu)\) belongs to Feigin--Fuchs
   Class~\(\mathrm{III}_-\): for every \(i\geq1\), its \(i\)-th Jantzen layer
   (under \(\Delta\mapsto\Delta+x\) at fixed \(c_\star\)) is
   \(J^i(c_\star,\mu)=M(c_\star,\mu^+_{k+i})+M(c_\star,\mu^-_{k+i})\), with
   \(M(c_\star,\mu^+_{k+i})\cap M(c_\star,\mu^-_{k+i})
   =J^{i+1}(c_\star,\mu)\) and
   \(J^i(c_\star,\mu)/J^{i+1}(c_\star,\mu)
   \cong L(c_\star,\mu^+_{k+i})\oplus L(c_\star,\mu^-_{k+i})\).
 \item[(ii)] The two primitive singular vectors of \(M(c_\star,\mu)\) generate
   \(M(c_\star,\mu^+_{k+1})\) and \(M(c_\star,\mu^-_{k+1})\).  Each primitive
   arrow \(\mu\to\nu\) (\(\nu\in\{\mu^+_{k+1},\mu^-_{k+1}\}\)) has a unique
   positive Kac label \((R,S)\in\mathbb{Z}_{\geq1}^2\) with
   \(h_{R,S}(b_\star)=\mu\) and \(h_{R,-S}(b_\star)=\nu\), determined by
   \(p'R+pS=n_\nu\) and \(p'R-pS=\pm n_\mu\) (with \(n_\mu:=n^\epsilon_k\)).
 \item[(iii)] Every primitive BGG arrow \((R,S)\) satisfies:
   \begin{align}
    &R\equiv r\text{ or }p-r\pmod p, \qquad S\equiv s\text{ or }p'-s\pmod{p'},
      \label{eq:RS-mod}\\
    &R\geq\min(r,p-r), \qquad S\geq\min(s,p'-s),
      \label{eq:RS-lower}\\
    &\min(R/p,\;S/p')<1.
      \label{eq:RS-small}
   \end{align}
   In particular, no integer pair \((qp,qp')\) with \(q\neq0\) lies in
   \([1-R,R]\times[1-S,S]\), so the first-layer Jantzen norm
   \(\gamma_{R,S}(b_\star):=\left.\partial_\Delta\right|_{\Delta=\mu}
   B_{c_\star,\Delta}(\chi_{R,S}(b_\star),\chi_{R,S}(b_\star))\) is nonzero.
 \item[(iv)] The two primitive arrows \((R_1,S_1)\) and \((R_2,S_2)\) out of
   \(\mu\) have distinct levels \(R_1S_1\neq R_2S_2\).  Setting
   \(m_i:=p'R_i-pS_i\in\{\pm n_\mu\}\) and \(n_i:=p'R_i+pS_i>0\) (\(i=1,2\)),
   we have
   \begin{equation}\label{eq:transverse-slope}
    \left.\frac{d}{dt}\right|_{t=t_\star}
    \bigl(h_{R_1,S_1}(t)-h_{R_2,S_2}(t)\bigr)
    =\frac{m_2n_2-m_1n_1}{4p'^2}\neq0,
    \qquad t_\star=-\frac{p'}{p}.
   \end{equation}
   Moreover, \(\max(\ell^+_1,\ell^-_1)<\min(\ell^+_2,\ell^-_2)\).
\end{enumerate}
\end{lemma}

\begin{proof}
Part~(i) follows from \cite[Theorem~A(a),(b)(iii),(c) and
Theorem~B(e)]{Astashkevich1997} (and \cite[Proposition~5.4, Figure~5.4, and
Theorem~6.3]{IoharaKoga2011}): since \(n_\mu\equiv\pm n_0\) or \(\pm n_1\pmod{2pp'}\)
is divisible by neither \(p\) nor \(p'\), the line \(p'x-py=n_\mu\) meets
neither axis at an integer point, placing \(M(c_\star,\mu)\) in
Class~\(\mathrm{III}_-\), and the fixed-\(c\) line \(c=c_\star\) is
transverse to every Kac curve through \((c_\star,\mu)\) because
\(dc/dt=6(1-t^{-2})\neq0\) at \(t_\star=-p'/p<0\).
For part~(ii), \(h_{R,S}(b_\star)=\mu\) and \(h_{R,-S}(b_\star)=\nu\) are
equivalent by \eqref{eq:mm-weight} to \(p'R+pS=n_\nu\) and
\(p'R-pS=\sigma n_\mu\) for some \(\sigma\in\{+1,-1\}\), i.e.\
\(2p'R=n_\nu+\sigma n_\mu\) and \(2pS=n_\nu-\sigma n_\mu\).  Because
\(p\nmid n_\mu\) and \(p'\nmid n_\mu\), at most one sign \(\sigma\) can make
\(n_\nu+\sigma n_\mu\) divisible by \(2p'\), giving uniqueness; the
consecutive-layer formulas in \eqref{eq:bgg-momenta} supply existence with
\(R,S\geq1\).

For part~(iii), reducing \(p'R-pS=\sigma n_\mu\equiv\pm(p'r\pm ps)\pmod{2pp'}\)
directly modulo \(p\) and modulo \(p'\) yields \(p'R\equiv\pm p'r\pmod p\) and
\(-pS\equiv\pm ps\pmod{p'}\).  Since \(\gcd(p,p')=1\) and \(R,S\geq1\), this
gives \eqref{eq:RS-mod} and \eqref{eq:RS-lower}.  Furthermore,
\eqref{eq:bgg-momenta} gives \(0<n_\nu-|n_\mu|<2pp'\) for every consecutive
pair \((\mu,\nu)\): indeed, for \(k=2j\) even,
\(n^+_{2j+1}-n^\pm_{2j}=n_1\mp n_0\in\{2ps,2p'r\}\subset(0,2pp')\) and
\(n^-_{2j+1}-n^\pm_{2j}=2pp'-n_1\mp n_0\in\{2p(p'-s),2p'(p-r)\}\subset(0,2pp')\);
for \(k=2j+1\) odd, the four differences are likewise
\(\{2p(p'-s),2p'(p-r),2ps,2p'r\}\subset(0,2pp')\).  Whichever sign
\(\sigma=\pm1\) satisfies \(\sigma n_\mu=-|n_\mu|\) or \(+|n_\mu|\), one of
\(2p'R=n_\nu-|n_\mu|\) or \(2pS=n_\nu-|n_\mu|\) is strictly less than \(2pp'\),
proving \eqref{eq:RS-small}.  By \cite[Proposition~6.2]{CompanionTorus} (which
derives the fixed-\(c\) derivative normalization from Zamolodchikov's
\(\alpha\)-derivative formula \cite[Section~6, equation~(6.1) and Section~7,
equation~(7.5)]{Zamolodchikov2004}, removing the top factor by dividing by
\(Rb+Sb^{-1}\)),
\(\gamma_{R,S}(b)^{-1}\propto A_{R,S}(b)
=\frac{1}{2}\prod_{(a,b')\in[1-R,R]\times[1-S,S]\setminus\{(0,0),(R,S)\}}(ab+b'b^{-1})^{-1}\);
since \eqref{eq:RS-small} precludes \(a=qp,b'=qp'\) (\(q\neq0\)), every factor
is nonzero at \(b_\star\), so \(\gamma_{R,S}(b_\star)\neq0\).

For part~(iv), \(R_1S_1=R_2S_2\) would mean \(n^+_{k+1}=n^-_{k+1}\), which
requires \(n_0=0\) (for \(k+1\) even) or \(n_1=pp'\) (for \(k+1\) odd), both
impossible since \(0<r<p\), \(0<s<p'\), and \(\gcd(p,p')=1\).  Differentiating
\(h_{R,S}(t)=\frac{1-R^2}{4}t+\frac{1-RS}{2}+\frac{1-S^2}{4t}\) gives
\(h'_{R,S}(t)=\frac{1-R^2}{4}-\frac{1-S^2}{4t^2}\), so at \(t_\star=-p'/p\),
\[
 4p'^2\left.\frac{d}{dt}\right|_{t=t_\star}
 \bigl(h_{R_1,S_1}(t)-h_{R_2,S_2}(t)\bigr)
 =p'^2(R_2^2-R_1^2)-p^2(S_2^2-S_1^2)
 =m_2n_2-m_1n_1.
\]
Since \(m_1,m_2\in\{\pm n_\mu\}\) with \(n_\mu\neq0\) and \(n_1,n_2\) are
distinct positive integers, \(m_2n_2-m_1n_1\) equals either
\(\pm n_\mu(n_2-n_1)\neq0\) (if \(m_1=m_2\)) or \(\pm n_\mu(n_2+n_1)\neq0\)
(if \(m_1=-m_2\)), proving \eqref{eq:transverse-slope}.  Finally,
\(\ell^+_1=rs<\min(p'r,ps)<pp'-|p'r-ps|=\min(\ell^+_2,\ell^-_2)\), and
similarly \(\ell^-_1=(p-r)(p'-s)<\min(p'(p-r),p(p'-s))<\min(\ell^+_2,\ell^-_2)\).
\end{proof}

\section{Fusion-Descent, the Jantzen Vertex Filtration, and Regularity}\label{sec:descent}

We first verify the injectivity and exact intersection property of the fusion
set \(\mathcal{F}_{p,p'}(r,s)\).

\begin{lemma}[Exact intersection of primitive fusion and residue roots]\label{lem:fusion-injective}
Let \(p,p'\geq2\) be coprime and \(0<r<p\), \(0<s<p'\).
\begin{enumerate}
 \item[(i)] The common roots \(d\in\mathbb{C}\) of
   \(P_{r,s}(b_\star;h_{r,s}(b_\star),d)=0\) and
   \(P_{p-r,p'-s}(b_\star;h_{r,s}(b_\star),d)=0\) are precisely the elements of
   \(\mathcal{F}_{p,p'}(r,s)\) in \eqref{eq:fusion-set}.
 \item[(ii)] The map \((k,m)\mapsto h_{2k+1,2m+1}(b_\star)\) is injective on
   \(0\leq k<\min(r,p-r)\), \(0\leq m<\min(s,p'-s)\).
 \item[(iii)] The common roots \(d\in\mathbb{C}\) of
   \(\tau_{r,s}(b_\star,d)=0\) and \(\tau_{p-r,p'-s}(b_\star,d)=0\) are
   precisely the elements of \(\mathcal{F}^\pm_{p,p'}(r,s)\) in
   \eqref{eq:fusion-set-pm}.
\end{enumerate}
\end{lemma}

\begin{proof}
By \eqref{eq:P-RS-roots}, a common root \(d\) in~(i) has representations
\(d=h_{u_1,v_1}(b_\star)=h_{u_2,v_2}(b_\star)\) with odd integers
\(1\leq u_1\leq 2r-1\), \(1\leq v_1\leq 2s-1\) and
\(1\leq u_2\leq 2(p-r)-1\), \(1\leq v_2\leq 2(p'-s)-1\).  By
\eqref{eq:mm-weight}, \(p'u_1-pv_1=\pm(p'u_2-pv_2)\).
If the sign is \(+\), then \(p'(u_1-u_2)=p(v_1-v_2)\); since \(\gcd(p,p')=1\)
and \(u_1-u_2,v_1-v_2\) are even, \(2p\mid(u_1-u_2)\).  Because
\(|u_1-u_2|\leq\max(2r-2,2(p-r)-2)<2p\), we get \(u_1=u_2\) and \(v_1=v_2\),
so \((u_1,v_1)=(2k+1,2m+1)\) lies in the intersection box
\eqref{eq:fusion-set}.
If the sign is \(-\), then \(p'(u_1+u_2)=p(v_1+v_2)\); since \(u_1+u_2\) and
\(v_1+v_2\) are positive even integers and \(\gcd(p,p')=1\), we must have
\(u_1+u_2\geq 2p\).  However,
\(u_1+u_2\leq (2r-1)+(2(p-r)-1)=2p-2<2p\), a contradiction.  This proves~(i).
For~(ii), if two pairs \((u_1,v_1),(u_2,v_2)\) in the intersection box have the
same weight, the same argument applies since
\(u_1+u_2\leq 4\min(r,p-r)-2\leq 2p-2<2p\) and \(|u_1-u_2|<p\).

For~(iii), recall from \cite[Proposition~6.2]{CompanionTorus} (or from the
factorization
\[
 \tau_{R,S}(b,d)=C_{R,S}(b)^2\,P_{R,S}\bigl(b;h_{R,S}(b),d\bigr)\,
                 P_{R,S}\bigl(b;h_{R,-S}(b),d\bigr)
\]
proved in Section~\ref{sec:cocycle}) that the roots of
\(\tau_{R,S}(b_\star,d)=0\) are \(d=h_{u,\pm v}(b_\star)\) with
\(u\in\{1,3,\ldots,2R-1\}\), \(v\in\{1,3,\ldots,2S-1\}\).  A common root of
\(\tau_{r,s}\) and \(\tau_{p-r,p'-s}\) therefore has
\(p'u_1-\epsilon_1pv_1=\pm(p'u_2-\epsilon_2pv_2)\) with
\(\epsilon_1,\epsilon_2\in\{+1,-1\}\).  If \(\epsilon_1=\epsilon_2=+1\),
part~(i) gives \(d\in\mathcal{F}_{p,p'}(r,s)\).  If
\(\epsilon_1=\epsilon_2=-1\), the minus sign is impossible because both sides
are positive, while the plus sign gives \(p'(u_1-u_2)=p(v_2-v_1)\), forcing
\(u_1=u_2\) and \(v_1=v_2\) as above, so \(d\in\mathcal{F}^\pm_{p,p'}(r,s)\).
If \(\epsilon_1=-\epsilon_2\) (say \(\epsilon_1=+1,\epsilon_2=-1\)), then
either \(p'(u_1-u_2)=p(v_1+v_2)\) (requiring \(2p'\mid(v_1+v_2)\), impossible
since \(0<v_1+v_2\leq 2p'-2<2p'\)) or \(p'(u_1+u_2)=p(v_1-v_2)\) (requiring
\(2p\mid(u_1+u_2)\), impossible since \(0<u_1+u_2\leq 2p-2<2p\)).
\end{proof}

Next we construct the Jantzen vertex filtration at any fusion weight
\(d\in\mathcal{F}_{p,p'}(r,s)\).  Fix any \(\mu\in\mathcal{O}_{p,p'}(r,s)\)
(say at layer \(k_0\geq0\)) and identify \(M(c_\star,\mu+x)_N[[x]]\) with
\(\mathbb{C}[[x]]^{p(N)}\) via the PBW basis.  For \(i\geq0\), define the
\(x\)-adic Jantzen lattice
\begin{equation}\label{eq:x-adic-jantzen}
 \mathcal{J}^{(i)}_N
 :=\bigl\{\,a(x)\in M(c_\star,\mu+x)_N[[x]]
   : B_x(a(x),v)\in x^i\mathbb{C}[[x]]\;\;\forall\,v\in M(c_\star,\mu+x)_N[[x]]\,\bigr\},
\end{equation}
where \(B_x:=B_{c_\star,\mu+x}\) and \(\rho_x:=\rho_{c_\star,\mu+x,d}\), and
set \(\mathcal{J}^{(i)}=\bigoplus_{N\geq0}\mathcal{J}^{(i)}_N\).  Since
\(B_x(L_ma(x),v)=B_x(a(x),L_{-m}v)\), each \(\mathcal{J}^{(i)}\) is a
\(\Vir\)-submodule of \(M(c_\star,\mu+x)[[x]]\).  By definition of the Jantzen
filtration \cite[Definition~5.2]{CompanionTorus}, evaluating at \(x=0\) maps
\(\mathcal{J}^{(i)}\) onto \(J^i(c_\star,\mu)\), and clearly
\begin{equation}\label{eq:x-shift-lattice}
 \mathcal{J}^{(i)}\cap x\,M(c_\star,\mu+x)[[x]] = x\,\mathcal{J}^{(i-1)}
 \qquad (i\geq1).
\end{equation}

\begin{proposition}[Jantzen vertex filtration and all-order pole cancellation]\label{prop:jantzen-vertex}
Let \(\mu\in\mathcal{O}_{p,p'}(r,s)\) (at layer \(k_0\geq0\)) and let
\(d=h_{2k+1,2m+1}(b_\star)\in\mathcal{F}_{p,p'}(r,s)\).
\begin{enumerate}
 \item[(i)] For every \(i\geq0\),
   \(\rho_x(\mathcal{J}^{(i)},\mathcal{J}^{(i)})\subset x^i\mathbb{C}[[x]]\),
   and the \(i\)-th Jantzen vertex form
   \(\rho^{(i)}(a(0),b(0)):=[x^i]\rho_x(a(x),b(x))\) is well-defined on
   \(J^i(c_\star,\mu)\times J^i(c_\star,\mu)\), satisfies the exact \(V_d(1)\)
   Ward identity \eqref{eq:two-weight-ward} (with \(\Delta_1=\Delta_2=\mu\)),
   and satisfies \(J^{i+1}(c_\star,\mu)\subset\rad\rho^{(i)}\).
 \item[(ii)] For every \(N\geq0\), \(F_N(c_\star,\mu+x,d)\) is holomorphic at
   \(x=0\) (all Jantzen poles cancel to order zero), and its value at \(x=0\)
   decomposes over the Jantzen layers as
   \begin{equation}\label{eq:jantzen-trace-decomp}
    F_N(c_\star,\mu,d)
    =T^{\mathrm{irr}}_N(c_\star,\mu,d)
     +\sum_{i=1}^{j_N(\mu)}\sum_{\pm}
       \frac{\sigma^\pm_i(\mu)}{\gamma^\pm_i(\mu)}\,
       T^{\mathrm{irr}}_{N-(\mu^\pm_{k_0+i}-\mu)}(c_\star,\mu^\pm_{k_0+i},d),
   \end{equation}
   where \(\gamma^\pm_i(\mu)=B^{(i)}(w^\pm_{k_0+i},w^\pm_{k_0+i})\neq0\) and
   \(\sigma^\pm_i(\mu)=\rho^{(i)}(w^\pm_{k_0+i},w^\pm_{k_0+i})\).
\end{enumerate}
\end{proposition}

\begin{proof}
We first record a general Ward-induction fact: if a symmetric bilinear form
\(\eta\) on a Verma module \(M(c_\star,\nu)\) satisfies the \(V_d(1)\) Ward
identity \eqref{eq:two-weight-ward} and \(w\in M(c_\star,\nu)_\ell\) is a
singular vector with \(\eta(w,v_\nu)=0\), then
\(U(\Vir_-)w\subset\rad\eta\).  Indeed, for \(n>0\) and homogeneous \(u\),
\(L_nw=0\) and \eqref{eq:two-weight-ward} give
\(\eta(w,L_{-n}u)=-(\ell-|u|-nd)\eta(w,u)\), so induction on \(|u|\) yields
\(\eta(w,M(c_\star,\nu))=0\); a second induction on PBW length using
\(\eta(L_{-n}a,u)=\eta(a,L_nu)+(|a|-|u|+nd)\eta(a,u)\) gives
\(\eta(U(\Vir_-)w,M(c_\star,\nu))=0\), and symmetry yields the other slot.

We prove~(i) by induction on \(i\geq0\).  For \(i=0\),
\(\rho^{(0)}=\rho_{c_\star,\mu,d}\).  By Lemma~\ref{lem:bgg-arrows}(i)--(iii),
\(J^1(c_\star,\mu)=U(\Vir_-)w^+_{k_0+1}+U(\Vir_-)w^-_{k_0+1}\), where each
primitive singular vector has Kac label \((R,S)\) with
\(R\geq\min(r,p-r)>k\) and \(S\geq\min(s,p'-s)>m\).  By
Proposition~\ref{prop:fusion-poly}(ii),
\(P_{R,S}(b_\star;\mu,h_{2k+1,2m+1}(b_\star))=0\), so
\(\rho^{(0)}(w^\pm_{k_0+1},v_\mu)=0\).  The Ward-induction fact above gives
\(J^1(c_\star,\mu)\subset\rad\rho^{(0)}\).

Now assume~(i) holds for \(i-1\geq0\).  If \(a(x),b(x)\in\mathcal{J}^{(i)}\),
then \(a(x),b(x)\in\mathcal{J}^{(i-1)}\) and
\(a(0),b(0)\in J^i\subset\rad\rho^{(i-1)}\), so
\([x^{i-1}]\rho_x(a(x),b(x))=\rho^{(i-1)}(a(0),b(0))=0\), whence
\(\rho_x(a(x),b(x))\in x^i\mathbb{C}[[x]]\).  If \(a_1(x),a_2(x)\in\mathcal{J}^{(i)}\)
have \(a_1(0)=a_2(0)\), then \(a_1(x)-a_2(x)=xw(x)\) with
\(w(x)\in\mathcal{J}^{(i-1)}\) by \eqref{eq:x-shift-lattice}; since
\(b(0)\in J^i\subset\rad\rho^{(i-1)}\), we get
\([x^{i-1}]\rho_x(w(x),b(x))=\rho^{(i-1)}(w(0),b(0))=0\), so
\(\rho_x(xw(x),b(x))\in x^{i+1}\mathbb{C}[[x]]\).  By symmetry, \(\rho^{(i)}\)
is independent of the lifts in \(\mathcal{J}^{(i)}\).  For homogeneous
\(a(x),b(x)\in\mathcal{J}^{(i)}\) of levels \(N_a,N_b\), both \(L_{-n}a(x)\)
and \(L_nb(x)\) lie in the \(\Vir\)-submodule \(\mathcal{J}^{(i)}\), and in
\eqref{eq:two-weight-ward} the weight difference
\((\mu+x+N_a)-(\mu+x+N_b)=N_a-N_b\) is independent of \(x\); taking \([x^i]\)
gives the exact \(V_d(1)\) Ward identity for \(\rho^{(i)}\) on \(J^i\).

It remains to show \(J^{i+1}\subset\rad\rho^{(i)}\) for \(i\geq1\).  By
Lemma~\ref{lem:bgg-arrows}(i), \(J^i=M(\mu^+_{k_0+i})+M(\mu^-_{k_0+i})\) and
\(J^{i+1}=M(\mu^+_{k_0+i})\cap M(\mu^-_{k_0+i})
=U(\Vir_-)w^+_{k_0+i+1}+U(\Vir_-)w^-_{k_0+i+1}\).  On each Verma submodule \(M(\mu^\epsilon_{k_0+i})\subset J^i\)
(\(\epsilon\in\{+,-\}\)), the Ward identity and
Lemma~\ref{lem:two-weight-ward} force the restriction of \(\rho^{(i)}\) to
\(M(\mu^\epsilon_{k_0+i})\times M(\mu^\epsilon_{k_0+i})\) to equal
\(\sigma^\epsilon_i(\mu)\,\rho_{c_\star,\mu^\epsilon_{k_0+i},d}\).  Since each
\(w^\eta_{k_0+i+1}\) (\(\eta\in\{+,-\}\)) is a primitive singular vector in
\emph{both} \(M(\mu^+_{k_0+i})\) and \(M(\mu^-_{k_0+i})\) with Kac label
\((R,S)\) satisfying \(R>k\) and \(S>m\) by \eqref{eq:RS-lower}, we obtain
\(\rho^{(i)}(w^\eta_{k_0+i+1},w^+_{k_0+i})=0\) (viewing
\(w^\eta_{k_0+i+1}\in M(\mu^+_{k_0+i})\)) and
\(\rho^{(i)}(w^\eta_{k_0+i+1},w^-_{k_0+i})=0\) (viewing
\(w^\eta_{k_0+i+1}\in M(\mu^-_{k_0+i})\)).  Because
\(L_nw^\eta_{k_0+i+1}=0\) for all \(n>0\), Ward induction on PBW descendants of
\(w^+_{k_0+i}\) and \(w^-_{k_0+i}\) gives
\(\rho^{(i)}(w^\eta_{k_0+i+1},J^i)=0\), and a second Ward induction gives
\(\rho^{(i)}(J^{i+1},J^i)=0\), completing the induction for~(i).

For~(ii), fix \(N\geq0\).  Over the discrete valuation ring \(\mathbb{C}[[x]]\)
(of characteristic \(0\)), every nondegenerate symmetric matrix \(G_N(c_\star,\mu+x)\)
is congruent over \(\mathrm{GL}_{p(N)}(\mathbb{C}[[x]])\) to a diagonal matrix:
there exists \(U(x)\in\mathrm{GL}_{p(N)}(\mathbb{C}[[x]])\) with columns
\(v_1(x),\ldots,v_{p(N)}(x)\) such that
\(B_x(v_a(x),v_b(x))=\delta_{ab}x^{e_a}u_a(x)\), where \(e_a\geq0\) are the
local Smith exponents and \(u_a(0)\neq0\).  Thus
\(v_a(x)\in\mathcal{J}^{(e_a)}_N\), and for each \(i\geq0\) the vectors
\(\{v_a(0):e_a=i\}\) project to a basis of \(J^i_N/J^{i+1}_N\) diagonalizing
the nondegenerate \(i\)-th Jantzen form \(\bar B^{(i)}_N\) with diagonal
entries \(u_a(0)\).  Computing the trace in this basis yields
\[
 F_N(c_\star,\mu+x,d)
 =\Tr\bigl((U^{\mathsf{T}}G_NU)^{-1}(U^{\mathsf{T}}\rho_NU)\bigr)
 =\sum_{a=1}^{p(N)}\frac{\rho_x(v_a(x),v_a(x))}{x^{e_a}u_a(x)}.
\]
By part~(i), \(\rho_x(v_a(x),v_a(x))=x^{e_a}\rho^{(e_a)}(v_a(0),v_a(0))+O(x^{e_a+1})\).
Every factor \(x^{e_a}\) cancels, so \(F_N(c_\star,\mu+x,d)\) is regular at
\(x=0\) with value
\[
 F_N(c_\star,\mu,d)
 =\sum_{i=0}^{j_N(\mu)}\Tr_{J^i_N/J^{i+1}_N}
   \bigl((\bar B^{(i)}_N)^{-1}\bar\rho^{(i)}_N\bigr).
\]
For \(i=0\), \(J^0/J^1=L(c_\star,\mu)\) gives \(T^{\mathrm{irr}}_N(c_\star,\mu,d)\).
For \(i\geq1\), \(J^i/J^{i+1}\cong L(c_\star,\mu^+_{k_0+i})\oplus
L(c_\star,\mu^-_{k_0+i})\) is a direct sum of non-isomorphic irreducibles,
which are orthogonal under the contravariant form \(\bar B^{(i)}\); hence
\((\bar B^{(i)}_N)^{-1}\) is block-diagonal, so off-diagonal blocks of
\(\bar\rho^{(i)}_N\) drop out of the trace.  On each summand
\(L(c_\star,\mu^\pm_{k_0+i})\), uniqueness of the contravariant and primary
forms gives \(\bar B^{(i)}=\gamma^\pm_i(\mu)\bar B_{L(\mu^\pm_{k_0+i})}\)
(with \(\gamma^\pm_i(\mu)\neq0\) because \(\bar B^{(i)}\) is nondegenerate) and
\(\bar\rho^{(i)}=\sigma^\pm_i(\mu)\bar\rho_{L(\mu^\pm_{k_0+i}),d}\), yielding
\eqref{eq:jantzen-trace-decomp}.
\end{proof}

We also record the joint regularity in \((b,\Delta)\), which will be needed in
Section~\ref{sec:telescoping}.

\begin{lemma}[Joint regularity in \((b,\Delta)\)]\label{lem:joint-holomorphicity}
Let \(\lambda=h_{r,s}(b_\star)\) be an interior Kac weight, let
\(d_\star=h_{2k+1,\epsilon(2m+1)}(b_\star)\in\mathcal{F}^\pm_{p,p'}(r,s)\)
(\(\epsilon\in\{+,-\}\)), and set \(d(b):=h_{2k+1,\epsilon(2m+1)}(b)\).  Then
for every \(N\geq0\) and every \(\mu\in\mathcal{O}_{p,p'}(r,s)\), the rational
function \(F_N(c(b),\Delta,d(b))\in\mathbb{Q}(b,\Delta)\) is regular
(holomorphic) at \((b_\star,\mu)\in\mathbb{C}^2\).
\end{lemma}

\begin{proof}
Work in the Laurent-polynomial ring \(\mathbb{Q}[b,b^{-1},\Delta]\).  By the
Kac determinant formula,
\[
 \det G_N(c(b),\Delta)
 =C_N\prod_{RS\leq N}\bigl(\Delta-h_{R,S}(b)\bigr)^{p(N-RS)}
\]
with \(C_N\in\mathbb{Q}^\times\) independent of \((b,\Delta)\), so
\(\det G_N(c(b),\Delta)\) has unit \(\Delta\)-leading coefficient and no
vertical prime factor in \(\mathbb{Q}[b,b^{-1}]\).  Since the entries of
\(G_N(c(b),\Delta)\) and \(\rho_N(c(b),\Delta,d(b))\) lie in
\(\mathbb{Q}[b,b^{-1},\Delta]\), the trace \(F_N(c(b),\Delta,d(b))\) lies in
the localization of \(\mathbb{Q}[b,b^{-1},\Delta]\) at the Kac factors
\(\Delta-h_{R,S}(b)\) (\(RS\leq N\)).

At generic \(b\), the Poghossian--HJS recursion
\cite[Proposition~6.4]{CompanionTorus} writes \(H_M(c(b),\Delta,d(b))\) as a
sum of simple poles at \(\Delta=h_{R,S}(b)\) (\(RS\leq M\)) with residues
\(R_{R,S}(b,d(b))=A_{R,S}(b)\tau_{R,S}(b,d(b))\)
\cite[Proposition~6.2]{CompanionTorus}.  Since the explicit product for
\(\tau_{R,S}(b,d)\) vanishes at \(d=h_{u,\pm v}(b)\) for all odd
\(1\leq u\leq 2R-1\) and \(1\leq v\leq 2S-1\), the residue \(R_{R,S}(b,d(b))\)
vanishes identically in \(b\) whenever \(R\geq k+1\) and \(S\geq m+1\).  Hence
over \(\mathbb{Q}(b)\), every inactive linear Kac factor \(\Delta-h_{R,S}(b)\)
with \(R\geq k+1\) and \(S\geq m+1\) cancels from the denominator of
\(H_M(c(b),\Delta,d(b))\), and therefore also from
\(F_N(c(b),\Delta,d(b))=\sum_{j=0}^Np(j)\,H_{N-j}(c(b),\Delta,d(b))\), leaving
a denominator in \(\mathbb{Q}[b,b^{-1},\Delta]\) dividing
\(\prod_{(R,S)\in\mathcal{A}_N}(\Delta-h_{R,S}(b))\), where
\(\mathcal{A}_N:=\{(R,S):RS\leq N,\; R\leq k\text{ or }S\leq m\}\) (since
\(F_N\) already had no vertical prime factor in \(\mathbb{Q}[b,b^{-1}]\) and
each \(\Delta-h_{R,S}(b)\) is monic in \(\Delta\)).

Finally, for any \((R,S)\) with
\(h_{R,S}(b_\star)=\mu\in\mathcal{O}_{p,p'}(r,s)\), \eqref{eq:mm-weight} gives
\(p'R-pS=\pm n_\mu\equiv\pm(p'r\pm ps)\pmod{2pp'}\); reducing modulo \(p\) and
\(p'\) forces \(R\geq\min(r,p-r)\geq k+1\) and \(S\geq\min(s,p'-s)\geq m+1\),
so \((R,S)\notin\mathcal{A}_N\).  Thus every active factor
\(\Delta-h_{R,S}(b)\) is nonzero at \((b_\star,\mu)\), so the denominator is a
unit in the local ring at \((b_\star,\mu)\).
\end{proof}

\begin{proof}[Proof of Theorem~\ref{thm:descent-regularity}]
Injectivity on the box \eqref{eq:fusion-set} is
Lemma~\ref{lem:fusion-injective}(ii).

\emph{Part~(i):} By symmetry \eqref{eq:two-weight-sym},
\(\rho(c_\star,\lambda,d)\) descends to
\(L(c_\star,\lambda)\times L(c_\star,\lambda)\) if and only if
\(J^1(c_\star,\lambda)\subset\rad\rho(c_\star,\lambda,d)\).  By
Lemma~\ref{lem:bgg-arrows}(i) and the Ward-induction fact in the proof of
Proposition~\ref{prop:jantzen-vertex}, this holds if and only if
\(\rho(w^+_1,v_\lambda)=\rho(w^-_1,v_\lambda)=0\), which by
Proposition~\ref{prop:fusion-poly}(ii) is equivalent to
\(P_{r,s}(b_\star;\lambda,d)=0\) and \(P_{p-r,p'-s}(b_\star;\lambda,d)=0\), and
hence by Lemma~\ref{lem:fusion-injective}(i) to \(d\in\mathcal{F}_{p,p'}(r,s)\).

\emph{Part~(ii) and the final sentence:} If
\(d\in\mathcal{F}^\pm_{p,p'}(r,s)\), then Lemma~\ref{lem:joint-holomorphicity}
gives joint regularity of \(F_N(c(b),\Delta,h_{2k+1,\epsilon(2m+1)}(b))\) at
\((b_\star,\mu)\), and hence regularity of \(F_N(c_\star,\Delta,d)\) at
\(\Delta=\mu\), for every \(N\geq0\) and every
\(\mu\in\mathcal{O}_{p,p'}(r,s)\).  (When \(d\in\mathcal{F}_{p,p'}(r,s)\),
Proposition~\ref{prop:jantzen-vertex}(ii) gives an independent intrinsic proof
together with the Jantzen trace decomposition \eqref{eq:jantzen-trace-decomp}.)

Conversely, suppose \(d\notin\mathcal{F}^\pm_{p,p'}(r,s)\).  By
Lemma~\ref{lem:fusion-injective}(iii), at least one of
\(\tau^+_1(d)=\tau_{r,s}(b_\star,d)\) and
\(\tau^-_1(d)=\tau_{p-r,p'-s}(b_\star,d)\) is nonzero.
Let \(\ell_{\min}:=\min(\ell^+_1,\ell^-_1)\) and
\(\ell_{\max}:=\max(\ell^+_1,\ell^-_1)\), with corresponding branch labels
\(\mathrm{small}\) and \(\mathrm{big}\).  By Lemma~\ref{lem:bgg-arrows}(iv),
\(\ell_{\min}<\ell_{\max}<\min(\ell^+_2,\ell^-_2)\), so both levels lie in
the first Jantzen window \(j_N(\lambda)=1\).  By the top-layer residue formula
\cite[Theorem~5.9]{CompanionTorus}:
if \(\tau^{\mathrm{small}}_1(d)\neq0\), then
\(\Res_{\Delta=\lambda}F_{\ell_{\min}}(c_\star,\Delta,d)
=\tau^{\mathrm{small}}_1(d)/\gamma^{\mathrm{small}}_1\neq0\);
if \(\tau^{\mathrm{small}}_1(d)=0\) and \(\tau^{\mathrm{big}}_1(d)\neq0\), then
\[
 \Res_{\Delta=\lambda}F_{\ell_{\max}}(c_\star,\Delta,d)
 =\frac{\tau^{\mathrm{small}}_1(d)}{\gamma^{\mathrm{small}}_1}
  F_{\ell_{\max}-\ell_{\min}}(c_\star,\mu^{\mathrm{small}}_1,d)
 +\frac{\tau^{\mathrm{big}}_1(d)}{\gamma^{\mathrm{big}}_1}
 =\frac{\tau^{\mathrm{big}}_1(d)}{\gamma^{\mathrm{big}}_1}\neq0
\]
(noting that \(F_{\ell_{\max}-\ell_{\min}}(c_\star,\mu^{\mathrm{small}}_1,d)\)
is finite since \(\det G_M(c_\star,\mu^{\mathrm{small}}_1)\neq0\) for
\(M<\min(\ell^+_2,\ell^-_2)-\ell_{\min}\), by
\cite[Theorem~5.9]{CompanionTorus}).  Thus \(F_N(c_\star,\Delta,d)\) has a
simple pole at \(\Delta=\lambda\) for \(N=\ell_{\min}\) or \(N=\ell_{\max}\).
\end{proof}

\section{First-Layer Vertex Ratio and the BGG Cocycle}\label{sec:cocycle}

We now prove Theorem~\ref{thm:first-layer-ratio} and establish the properties
of the BGG cocycle at minimal-model charges.

\begin{proof}[Proof of Theorem~\ref{thm:first-layer-ratio}]
Fix generic \(b\in\mathbb{C}^\times\), \(t=b^2\), \(R,S\geq1\), and
\(0\leq k\leq R-1\), \(0\leq m\leq S-1\).  Let \(\mathcal{O}_{R,S}\in U(\Vir_-)_{RS}\)
be the constant-coefficient PBW operator such that
\(\chi_{R,S}(b)=\mathcal{O}_{R,S}v_{h_{R,S}(b)}\).  For any \(d\in\mathbb{C}\),
consider the two-variable polynomial
\[
 f_d(x_1,x_2)
 :=\rho_{h_{R,S}(b)+x_1,\;h_{R,S}(b)+x_2;\;d}
   \bigl(\mathcal{O}_{R,S}v_{h_{R,S}(b)+x_1},\;
         \mathcal{O}_{R,S}v_{h_{R,S}(b)+x_2}\bigr).
\]
By Lemma~\ref{lem:two-weight-ward}, \(f_d(x_1,x_2)=f_d(x_2,x_1)\).

Set \(x_2=0\).  Then \(\chi_{R,S}(b)=\mathcal{O}_{R,S}v_{h_{R,S}(b)}\) is a
singular vector of weight \(h_{R,-S}(b)=h_{R,S}(b)+RS\) in
\(M(c(b),h_{R,S}(b))\), generating a Verma submodule isomorphic to
\(M(c(b),h_{R,-S}(b))\).  Consequently, the map
\((u,w)\mapsto\rho_{h_{R,S}(b)+x_1,\,h_{R,S}(b);\,d}(u,w)\) restricted to
\(M(c(b),h_{R,S}(b)+x_1)\times U(\Vir_-)\chi_{R,S}(b)\) satisfies the exact
\(V_d(1)\) Ward identity \eqref{eq:two-weight-ward} between
\(M(c(b),h_{R,S}(b)+x_1)\) and \(M(c(b),h_{R,-S}(b))\).  By uniqueness in
Lemma~\ref{lem:two-weight-ward}, for every \(u\in M(c(b),h_{R,S}(b)+x_1)\),
\begin{multline*}
 \rho_{h_{R,S}(b)+x_1,\,h_{R,S}(b);\,d}\bigl(u,\chi_{R,S}(b)\bigr)\\
 =\rho_{h_{R,S}(b)+x_1,\,h_{R,S}(b);\,d}\bigl(v_{h_{R,S}(b)+x_1},\chi_{R,S}(b)\bigr)\,
  \rho_{h_{R,S}(b)+x_1,\,h_{R,-S}(b);\,d}\bigl(u,v_{h_{R,-S}(b)}\bigr).
\end{multline*}
Applying \eqref{eq:two-weight-sym} and Proposition~\ref{prop:fusion-poly}(ii)
to the first factor and substituting \(u=\mathcal{O}_{R,S}v_{h_{R,S}(b)+x_1}\)
yields
\[
 f_d(x_1,0)
 =C_{R,S}(b)\,P_{R,S}\bigl(b;h_{R,S}(b)+x_1,d\bigr)\,
  \rho_{h_{R,S}(b)+x_1,\,h_{R,-S}(b);\,d}
  \bigl(\mathcal{O}_{R,S}v_{h_{R,S}(b)+x_1},v_{h_{R,-S}(b)}\bigr).
\]
Now specialize to any \(d=h_{2k+1,2m+1}(b)\) with \(0\leq k<R\) and
\(0\leq m<S\) (including \(d=h_{1,1}(b)=0\) for \(k=m=0\)).  Then
\(P_{R,S}(b;h_{R,S}(b),d)=0\), while at \(x_1=0\) the second factor evaluates
by Proposition~\ref{prop:fusion-poly}(ii) to
\(C_{R,S}(b)P_{R,S}(b;h_{R,-S}(b),d)\).  By the product rule and the symmetry
\(f_d(x_1,x_2)=f_d(x_2,x_1)\),
\begin{equation}\label{eq:f-diag-deriv}
 \left.\frac{d}{dx}\right|_{x=0}f_d(x,x)
 =2\,\partial_{x_1}f_d(0,0)
 =2\,C_{R,S}(b)^2\,
  \left(\left.\partial_{\Delta_1}\right|_{\Delta_1=h_{R,S}(b)}
        P_{R,S}(b;\Delta_1,d)\right)
  P_{R,S}\bigl(b;h_{R,-S}(b),d\bigr).
\end{equation}
By \eqref{eq:d0-shapovalov}, \(f_0(x,x)=B_{c(b),h_{R,S}(b)+x}(\chi_{R,S},\chi_{R,S})\).
Dividing \eqref{eq:f-diag-deriv} for \(d=h_{2k+1,2m+1}(b)\) by
\eqref{eq:f-diag-deriv} for \(d=0\) cancels \(2C_{R,S}(b)^2\) identically:
\begin{equation}\label{eq:ratio-from-P}
 c_{R,S}(t;2k+1,2m+1)
 =\frac{\left.\partial_{\Delta_1}\right|_{\Delta_1=h_{R,S}(b)}
        P_{R,S}(b;\Delta_1,h_{2k+1,2m+1}(b))\cdot
        P_{R,S}(b;h_{R,-S}(b),h_{2k+1,2m+1}(b))}
       {\left.\partial_{\Delta_1}\right|_{\Delta_1=h_{R,S}(b)}
        P_{R,S}(b;\Delta_1,0)\cdot
        P_{R,S}(b;h_{R,-S}(b),0)}.
\end{equation}

To evaluate \eqref{eq:ratio-from-P}, rescale the momenta in
\eqref{eq:P-RS-def} by \(L_i:=b\lambda_i\) so that
\(\Delta_1=(Q^2-L_1^2/t)/4\) and \(\partial_{\Delta_1}=(-2t/L_1)\partial_{L_1}\);
the prefactor \(-2t/L_{\mathrm{in}}\) (where \(L_{\mathrm{in}}:=Rt+S\)) cancels
between numerator and denominator of \eqref{eq:ratio-from-P}.  Put
\(L_{\mathrm{out}}:=Rt-S\) and \(L_d:=(2k+1)t+(2m+1)\).  At generic \(t\), the
factorized product \(P_{R,S}(b;\Delta_1,d)\) in \eqref{eq:P-RS-def} has a
unique vanishing linear factor at \(L_1=L_{\mathrm{in}}\), namely
\(L_1-L_d+p_0t+q_0\) with \((p_0,q_0)=(2k+1-R,2m+1-S)\), whose
\(L_1\)-derivative is \(1\).  Grouping the remaining \(4RS-1\) linear factors
of \(\partial_{L_1}P_{R,S}|_{L_{\mathrm{in}}}\cdot P_{R,S}|_{L_{\mathrm{out}}}\)
(up to the \((k,m)\)-independent power of \(2b\)) into four quadrants indexed
by \((p',q')\in\{1-R,3-R,\ldots,R-1\}\times\{1-S,3-S,\ldots,S-1\}\):
\begin{align*}
 \tfrac{1}{2}(L_{\mathrm{in}}+L_d+p't+q')
 &=(a+k)t+(b'+m),
 \quad (a,b')\in[1,R]\times[1,S],\\
 -\tfrac{1}{2}(L_{\mathrm{in}}-L_d+p't+q')
 &=(a+k)t+(b'+m),
 \quad (a,b')\in[1-R,0]\times[1-S,0],\\
 \tfrac{1}{2}(L_{\mathrm{out}}+L_d+p't+q')
 &=(a+k)t+(b'+m),
 \quad (a,b')\in[1,R]\times[1-S,0],\\
 -\tfrac{1}{2}(L_{\mathrm{out}}-L_d+p't+q')
 &=(a+k)t+(b'+m),
 \quad (a,b')\in[1-R,0]\times[1,S],
\end{align*}
where the four bijections are \((a,b')=((R+1+p')/2,(S+1+q')/2)\),
\((-(R-1+p')/2,-(S-1+q')/2)\), \(((R+1+p')/2,(-S+1+q')/2)\), and
\((-(R-1+p')/2,(S+1-q')/2)\), respectively,
with the unique omitted factor at \((a,b')=(-k,-m)\) in the second line.  The
four quadrants partition \([1-R,R]\times[1-S,S]\setminus\{(-k,-m)\}\), and the
total sign \((-1)^{2RS-1}\) and power of \(2b\) are independent of \((k,m)\).
Dividing by the \(k=m=0\) case therefore gives \eqref{eq:box-formula}.

For the Gamma factorization \eqref{eq:gamma-potential}, cancel the common
intersection between the two boxes in \eqref{eq:box-formula} in two steps:
first shift \(b'\) by \(m\) at fixed \(a\in[1-R,R]\), then shift \(a\) by \(k\)
at fixed \(b'\in[1-S+m,S+m]\).  This writes \eqref{eq:box-formula} as
\begin{equation}\label{eq:two-step-box}
 c_{R,S}(t;2k+1,2m+1)
 =\left(\prod_{j=1}^m\prod_{a=1-R}^R\frac{at+S+j}{at-S+j}\right)
  \left(\prod_{i=1}^k\prod_{b'=1-S+m}^{S+m}\frac{(R+i)t+b'}{(-R+i)t+b'}\right).
\end{equation}
Set \(x_{\mathrm{in}}:=2P_{R,S}/b=R+S/t\), \(x_{\mathrm{out}}:=2P_{R,-S}/b=R-S/t\),
\(y_{\mathrm{in}}:=2P_{R,S}b=Rt+S\), and \(y_{\mathrm{out}}:=2P_{R,-S}b=Rt-S\).
Since \((1+x_{\mathrm{in}}+j/t)-(1-x_{\mathrm{out}}+j/t)=2R\in\mathbb{Z}_{>0}\),
\((1+x_{\mathrm{out}}+j/t)-(1-x_{\mathrm{in}}+j/t)=2R\in\mathbb{Z}_{>0}\),
\((1+y_{\mathrm{in}}+it+m)-(1+y_{\mathrm{out}}+it+m)=2S\in\mathbb{Z}_{>0}\),
and \((1-y_{\mathrm{out}}+it+m)-(1-y_{\mathrm{in}}+it+m)=2S\in\mathbb{Z}_{>0}\),
iterating \(\Gamma(z+1)=z\Gamma(z)\) gives
\begin{align*}
 \prod_{a=1-R}^R\frac{a+(S+j)/t}{a+(-S+j)/t}
 &=\frac{\Gamma(1+x_{\mathrm{in}}+j/t)}{\Gamma(1-x_{\mathrm{out}}+j/t)}
   \frac{\Gamma(1-x_{\mathrm{in}}+j/t)}{\Gamma(1+x_{\mathrm{out}}+j/t)},\\
 \prod_{b'=1-S+m}^{S+m}\frac{(R+i)t+b'}{(-R+i)t+b'}
 &=\frac{\Gamma(1+y_{\mathrm{in}}+it+m)}{\Gamma(1+y_{\mathrm{out}}+it+m)}
   \frac{\Gamma(1-y_{\mathrm{in}}+it+m)}{\Gamma(1-y_{\mathrm{out}}+it+m)}.
\end{align*}
Multiplying over \(j=1,\ldots,m\) and \(i=1,\ldots,k\) gives
\[
 c_{R,S}(t;2k+1,2m+1)
 =\frac{\Psi_{2k+1,2m+1}(P_{R,-S};b)}{\Psi_{2k+1,2m+1}(P_{R,S};b)}.
\]
\end{proof}

\begin{proposition}[Finiteness, nonvanishing, and path-independence of the BGG cocycle]\label{prop:bgg-cocycle}
Let \(b_\star^2=-p'/p\) (\(p,p'\geq2\) coprime), \(\lambda=h_{r,s}(b_\star)\)
interior, and \(d=h_{2k+1,2m+1}(b_\star)\in\mathcal{F}_{p,p'}(r,s)\).
\begin{enumerate}
 \item[(i)] For every primitive BGG arrow \(\mu\to\nu\) in
   \(\mathcal{O}_{p,p'}(r,s)\) with Kac label \((R,S)\), setting
   \(B_0:=[1-R,R]\times[1-S,S]\) and \(B_1:=[1-R+k,R+k]\times[1-S+m,S+m]\),
   every factor \(pb'-p'a\) with \((a,b')\in(B_1\setminus B_0)\cup(B_0\setminus B_1)\)
   is nonzero.  Consequently,
   \begin{equation}\label{eq:reduced-box-mm}
    c_{R,S}\Bigl(-\frac{p'}{p};2k+1,2m+1\Bigr)
    =\frac{\prod_{(a,b')\in B_1\setminus B_0}(pb'-p'a)}
          {\prod_{(a,b')\in B_0\setminus B_1}(pb'-p'a)}
    \;\in\;\mathbb{Q}^\times
   \end{equation}
   is finite and nonzero.
 \item[(ii)] For every diamond
   \(\mu\to\nu_1\to\xi\), \(\mu\to\nu_2\to\xi\) in the braided BGG diagram of
   \(\mathcal{O}_{p,p'}(r,s)\),
   \begin{equation}\label{eq:diamond-cocycle}
    c(\mu\to\nu_1)\,c(\nu_1\to\xi)
    =c(\mu\to\nu_2)\,c(\nu_2\to\xi),
   \end{equation}
   where \(c(\alpha\to\beta):=c_{R,S}(-p'/p;2k+1,2m+1)\) for the primitive
   arrow \(\alpha\to\beta\) of label \((R,S)\).  Hence the product of
   \(c(\alpha\to\beta)\) along any directed BGG path from \(\mu\) to \(\xi\)
   depends only on \((\mu,\xi)\), defining a canonical rational cocycle
   \(\Omega_{2k+1,2m+1}(\mu\to\xi)\in\mathbb{Q}^\times\).
\end{enumerate}
\end{proposition}

\begin{proof}
For~(i), since \(\gcd(p,p')=1\), \(pb'-p'a=0\) requires \(p\mid a\) and
\(p'\mid b'\).  If \((a,b')\in B_1\setminus B_0\), then either
\(a=R+i\) with \(1\leq i\leq k\) or \(b'=S+j\) with \(1\leq j\leq m\).
By \eqref{eq:RS-mod} and \(k<\min(r,p-r)\), if \(a=R+i\) (\(1\leq i\leq k\))
then \(a\equiv r+i\in\{r+1,\ldots,p-1\}\) (when \(R\equiv r\pmod p\)) or
\(a\equiv p-r+i\in\{p-r+1,\ldots,p-1\}\) (when \(R\equiv p-r\pmod p\)), so
\(p\nmid a\).  By the same check on \(b'=S+j\) (\(1\leq j\leq m\)) using
\(m<\min(s,p'-s)\), we have \(p'\nmid b'\).  Hence on \(B_1\setminus B_0\) the
two divisibilities \(p\mid a\) and \(p'\mid b'\) cannot both hold.  If
\((a,b')\in B_0\setminus B_1\), then either \(a=i-R\) (\(1\leq i\leq k\)) or
\(b'=j-S\) (\(1\leq j\leq m\)), and the same congruences give
\(a\equiv i-r\) or \(i-(p-r)\in\{1-p,\ldots,-1\}\) (so \(p\nmid a\)) or
\(p'\nmid b'\); again the two divisibilities cannot both hold.  Since
\(|B_1\setminus B_0|=|B_0\setminus B_1|\), the common powers of \(p^{-1}\) in
\(a(-p'/p)+b'=(pb'-p'a)/p\) cancel between numerator and denominator, giving
\eqref{eq:reduced-box-mm}.

For~(ii), set \(z_\mu:=n_\mu/p'\) for each \(\mu\in\mathcal{O}_{p,p'}(r,s)\)
(so \(2P_\mu/b_\star=\pm z_\mu\) and \(2P_\mu b_\star=\pm t_\star z_\mu=\mp n_\mu/p\)),
and evaluate the Gamma potential \eqref{eq:gamma-potential} directly at
\(b=b_\star\):
\[
 \Psi_{2k+1,2m+1}(P_\mu;b_\star)
 =\prod_{j=1}^m\frac{1}{\Gamma\bigl(1+\tfrac{n_\mu-jp}{p'}\bigr)
                         \Gamma\bigl(1+\tfrac{-n_\mu-jp}{p'}\bigr)}
  \prod_{i=1}^k\frac{1}{\Gamma\bigl(1+m+\tfrac{-n_\mu-ip'}{p}\bigr)
                         \Gamma\bigl(1+m+\tfrac{n_\mu-ip'}{p}\bigr)}.
\]
None of the Gamma arguments is an integer!  Indeed, by \eqref{eq:bgg-momenta},
\(n_\mu\equiv\pm ps\pmod{p'}\) and \(n_\mu\equiv\pm p'r\pmod{p}\).  Hence
\((n_\mu\pm jp)/p'\in\mathbb{Z}\) would require \(p'\mid(\pm s-j)\) for some
\(1\leq j\leq m<\min(s,p'-s)\); this is impossible because
\(1\leq s-j\leq s-1<p'\) and \(2\leq s+j\leq p'-1<p'\).  Similarly,
\((\pm n_\mu-ip')/p\in\mathbb{Z}\) would require \(p\mid(\pm r-i)\) for some
\(1\leq i\leq k<\min(r,p-r)\), which is impossible because
\(1\leq r-i\leq r-1<p\) and \(2\leq r+i\leq p-1<p\).  Since \(\Gamma\) has no
zeros and poles only at \(\mathbb{Z}_{\leq0}\), every factor in
\(\Psi_{2k+1,2m+1}(P_\mu;b_\star)\) is finite and nonzero, and
\(\Psi_{2k+1,2m+1}(P_\mu;b_\star)\in\mathbb{C}^\times\) depends only on \(\mu\)
(as \eqref{eq:gamma-potential} is even in \(P\)).

Now for any primitive BGG arrow \(\alpha\to\beta\) of label \((R,S)\) with
\(p'R-pS=\sigma n_\alpha\) (\(\sigma\in\{+1,-1\}\)) and \(p'R+pS=n_\beta\),
we have \(R+S/t_\star=\sigma n_\alpha/p'\), \(R-S/t_\star=n_\beta/p'\),
\(Rt_\star+S=-\sigma n_\alpha/p\), and \(Rt_\star-S=-n_\beta/p\).  Since
every Gamma value is finite and nonzero, the \(\Gamma(z+1)=z\Gamma(z)\)
telescoping in the proof of Theorem~\ref{thm:first-layer-ratio} holds pointwise
at \(t=t_\star=-p'/p\), giving
\begin{equation}\label{eq:arrow-is-psi-ratio}
 c_{R,S}\Bigl(-\frac{p'}{p};2k+1,2m+1\Bigr)
 =\frac{\Psi_{2k+1,2m+1}(P_\beta;b_\star)}{\Psi_{2k+1,2m+1}(P_\alpha;b_\star)}.
\end{equation}
Multiplying along the two sides of any BGG diamond immediately gives
\eqref{eq:diamond-cocycle}, and the path-independent product equals
\(\Omega_{2k+1,2m+1}(\mu\to\xi)
=\Psi_{2k+1,2m+1}(P_\xi;b_\star)/\Psi_{2k+1,2m+1}(P_\mu;b_\star)
\in\mathbb{Q}^\times\).
\end{proof}

\section{Two-Layer \texorpdfstring{\(t\)}{t}-Jantzen Telescoping and Generalized Rocha--Caridi}\label{sec:telescoping}

We now prove Theorem~\ref{thm:gen-rocha-caridi} and
Corollary~\ref{cor:collision-free}.

\begin{lemma}[Two-layer \(t\)-Jantzen identity on the generic degenerate quotient]\label{lem:degenerate-quotient-t}
Let \(b_\star^2=-p'/p\) (\(p,p'\geq2\) coprime), let
\(\lambda=h_{r,s}(b_\star)\) be an interior Kac weight, and let
\(d_\star=h_{u,v}(b_\star)\in\mathcal{F}_{p,p'}(r,s)\) with
\((u,v)=(2k+1,2m+1)\).  Let \(\mu\in\mathcal{O}_{p,p'}(r,s)\) have primitive
BGG arrows \(\mu\to\nu_1\) (label \((R_1,S_1)\)) and \(\mu\to\nu_2\) (label
\((R_2,S_2)\)).  Then as formal power series in \(q\),
\begin{equation}\label{eq:two-term-id}
 \mathcal{T}^{\mathrm{irr}}_{c_\star,\mu,d_\star}(q)
 +\Omega_{u,v}(\mu\to\nu_2)\,q^{\nu_2-\mu}\,
  \mathcal{T}^{\mathrm{irr}}_{c_\star,\nu_2,d_\star}(q)
 =\mathcal{F}^{\mathrm{Verma}}_{c_\star,\mu,d_\star}(q)
 -\Omega_{u,v}(\mu\to\nu_1)\,q^{\nu_1-\mu}\,
  \mathcal{F}^{\mathrm{Verma}}_{c_\star,\nu_1,d_\star}(q).
\end{equation}
\end{lemma}

\begin{proof}
Set \(t_0:=b_\star^2=-p'/p\), let \(\tau\) be a formal deformation parameter
with \(t(\tau):=t_0+\tau\) and \(b(\tau):=\sqrt{t(\tau)}\) near \(b_\star\),
and put \(\Delta_1(\tau):=h_{R_1,S_1}(b(\tau))\),
\(\Delta_2(\tau):=h_{R_2,S_2}(b(\tau))\), and
\(d(\tau):=h_{2k+1,2m+1}(b(\tau))\).  By Lemma~\ref{lem:bgg-arrows}(iii),
\(0\leq k<R_i\) and \(0\leq m<S_i\) for \(i=1,2\).

\emph{Step 1: Generic-\(\tau\) splitting.}
Fix generic \(\tau\neq0\) so that \(t(\tau)\) is irrational.  Then the module
\(M\bigl(c(b(\tau)),\Delta_1(\tau)\bigr)\) has a unique proper nonzero
submodule, generated by the singular vector \(\chi_{R_1,S_1}(b(\tau))\), which
has weight \(h_{R_1,-S_1}(b(\tau))=\Delta_1(\tau)+R_1S_1\).  Denote its
irreducible quotient by
\[
 V_{R_1,S_1}(b(\tau))
 :=M\bigl(c(b(\tau)),\Delta_1(\tau)\bigr)\big/
   U(\Vir_-)\chi_{R_1,S_1}(b(\tau)).
\]
Since \(0\leq k<R_1\) and \(0\leq m<S_1\), Proposition~\ref{prop:fusion-poly}(ii)
and the Ward-induction fact in Proposition~\ref{prop:jantzen-vertex} show that
\(U(\Vir_-)\chi_{R_1,S_1}(b(\tau))\) lies in the radical of both
\(B_{c(b(\tau)),\Delta_1(\tau)}\) and
\(\rho_{c(b(\tau)),\Delta_1(\tau),d(\tau)}\), so both forms descend to
\(V_{R_1,S_1}(b(\tau))\) with \(\bar B_{V,\tau}\) nondegenerate.  Applying
the \(\Delta\)-Jantzen decomposition (the exact two-layer analogue of
Proposition~\ref{prop:jantzen-vertex}(ii) for the generic Verma module
\(M(c(b(\tau)),\Delta_1(\tau))\)) together with
Theorem~\ref{thm:first-layer-ratio} gives, at every level \(N\geq0\) and
generic \(\tau\neq0\),
\begin{equation}\label{eq:generic-tau-split}
 \begin{aligned}
  F_N\bigl(c(b(\tau)),\Delta_1(\tau),d(\tau)\bigr)
  &=\Tr_{V_{R_1,S_1}(b(\tau))_N}\bigl(V_{d(\tau)}(1)\bigr)\\
  &\quad+c_{R_1,S_1}(t(\tau);u,v)\,
    F_{N-R_1S_1}\bigl(c(b(\tau)),h_{R_1,-S_1}(b(\tau)),d(\tau)\bigr).
 \end{aligned}
\end{equation}

\emph{Step 2: Free \(\mathbb{C}[[\tau]]\)-lattice and length-two specialization.}
By Proposition~\ref{prop:fusion-poly}(i), \(\chi_{R_1,S_1}(b(\tau))\) has
regular coefficients in \(\mathbb{C}[[\tau]]\) and specializes at \(\tau=0\) to
the nonzero singular vector \(\chi_{R_1,S_1}(b_\star)\in M(c_\star,\mu)_{R_1S_1}\).
Because \(M(c_\star,\mu)\) is free over \(U(\Vir_-)\), the PBW descendants
\(\{L_{-I}\chi_{R_1,S_1}(b_\star):|I|=N-R_1S_1\}\) are linearly independent
in \(M(c_\star,\mu)_N\).  Completing them to a basis of \(M(c_\star,\mu)_N\)
yields a free \(\mathbb{C}[[\tau]]\)-submodule complement, so the quotient
\(\mathcal{V}_N:=V_{R_1,S_1}(b(\tau))_N\) over \(\mathbb{C}[[\tau]]\) is free
of constant rank \(p(N)-p(N-R_1S_1)\), specializing at \(\tau=0\) to
\(V_0:=M(c_\star,\mu)/M(c_\star,\nu_1)\).

By Lemma~\ref{lem:bgg-arrows}(i), \(J^1(c_\star,\mu)=M(c_\star,\nu_1)+M(c_\star,\nu_2)\)
and \(M(c_\star,\nu_1)\cap M(c_\star,\nu_2)=J^2(c_\star,\mu)=\rad M(c_\star,\nu_2)\).
By the second isomorphism theorem,
\[
 \frac{J^1(c_\star,\mu)}{M(c_\star,\nu_1)}
 =\frac{M(c_\star,\nu_1)+M(c_\star,\nu_2)}{M(c_\star,\nu_1)}
 \cong\frac{M(c_\star,\nu_2)}{\rad M(c_\star,\nu_2)}
 =L(c_\star,\nu_2),
\]
and \(V_0/\bigl(J^1(c_\star,\mu)/M(c_\star,\nu_1)\bigr)
=M(c_\star,\mu)/J^1(c_\star,\mu)=L(c_\star,\mu)\).  Thus \(V_0\) has exact
length two with short exact sequence
\(0\to L(c_\star,\nu_2)\to V_0\to L(c_\star,\mu)\to0\).  Moreover, any
submodule \(W\subset V_0\) that maps onto the head \(L(c_\star,\mu)\) must
contain the unique highest-weight line of \(V_0\) (by the \(L_0\)-weight
grading) and hence equals the cyclic module \(V_0\); therefore the only
submodules of \(V_0\) are \(V_0\), \(L(c_\star,\nu_2)\), and \(0\).

\emph{Step 3: Exact depth-one \(\tau\)-Jantzen filtration on \(\mathcal{V}_N\).}
Define the \(\tau\)-Jantzen filtration
\(V_0=J_\tau^0\supseteq J_\tau^1\supseteq J_\tau^2\supseteq\cdots\) via
\(\mathcal{J}_\tau^{(i)}
:=\{a(\tau)\in\mathcal{V}:\bar B_{V,\tau}(a(\tau),v)\in\tau^i\mathbb{C}[[\tau]]
\;\forall\,v\in\mathcal{V}\}\).
Each \(J_\tau^i\) is a \(\Vir\)-submodule of \(V_0\), with
\(J_\tau^1=\rad\bar B_{V,0}=L(c_\star,\nu_2)\).  Since \(L(c_\star,\nu_2)\) is
irreducible, \(J_\tau^2\) is either \(L(c_\star,\nu_2)\) or \(0\).
Let \(w_2(\tau)\in\mathcal{V}_{R_2S_2}\) be the image of
\(\chi_{R_2,S_2}(b(\tau))\in M(c(b(\tau)),\Delta_1(\tau))_{R_2S_2}\)
(written with PBW coefficients polynomial in \(t(\tau),t(\tau)^{-1}\)).
Consider the two-variable regular functions
\[
 g_B(t,\Delta):=B_{c(b),\Delta}\bigl(\chi_{R_2,S_2}(b),\chi_{R_2,S_2}(b)\bigr),
 \qquad
 g_\rho(t,\Delta):=\rho_{c(b),\Delta,h_{u,v}(b)}\bigl(\chi_{R_2,S_2}(b),\chi_{R_2,S_2}(b)\bigr),
\]
where \(b=\sqrt{t}\).  Because \(\chi_{R_2,S_2}(b)\) is singular at
\(\Delta=h_{R_2,S_2}(t)\) and \(0\leq k<R_2,0\leq m<S_2\), both \(g_B\) and
\(g_\rho\) vanish identically along the curve \(\Delta=h_{R_2,S_2}(t)\).  Hence
\(g_B(t,\Delta)=(\Delta-h_{R_2,S_2}(t))\tilde g_B(t,\Delta)\) and
\(g_\rho(t,\Delta)=(\Delta-h_{R_2,S_2}(t))\tilde g_\rho(t,\Delta)\) for
functions \(\tilde g_B,\tilde g_\rho\) regular at \((t_0,\mu)\), with
\(\tilde g_B(t_0,\mu)=\gamma_{R_2,S_2}(b_\star)\neq0\) by
Lemma~\ref{lem:bgg-arrows}(iii).  Evaluating along \((t,\Delta)=(t_0+\tau,\Delta_1(\tau))\) and using
\eqref{eq:transverse-slope} gives
\(\Delta_1(\tau)-\Delta_2(\tau)=\alpha\tau+O(\tau^2)\) with
\(\alpha=(m_2n_2-m_1n_1)/(4p'^2)\neq0\), where \(m_i=p'R_i-pS_i\) and
\(n_i=p'R_i+pS_i\), so
\begin{equation}\label{eq:tau-deriv-chain}
 \begin{aligned}
  \bar B_{V,\tau}(w_2(\tau),w_2(\tau))
  &=g_B(t_0+\tau,\Delta_1(\tau))
   =\alpha\,\gamma_{R_2,S_2}(b_\star)\,\tau+O(\tau^2),\\
  \bar\rho_{V,\tau}(w_2(\tau),w_2(\tau))
  &=g_\rho(t_0+\tau,\Delta_1(\tau))
   =\alpha\,\tilde g_\rho(t_0,\mu)\,\tau+O(\tau^2).
 \end{aligned}
\end{equation}
For any lift \(\tilde w_2(\tau)\in\mathcal{V}_{R_2S_2}\) with
\(\tilde w_2(0)=w_2(0)\), we have \(\tilde w_2(\tau)-w_2(\tau)=\tau u(\tau)\),
and hence the difference
\(\bar B_{V,\tau}(\tilde w_2(\tau),\tilde w_2(\tau))
-\bar B_{V,\tau}(w_2(\tau),w_2(\tau))\) is \(O(\tau^2)\) because
\(w_2(0)\in J_\tau^1=\rad\bar B_{V,0}\).  Since
\(\alpha\,\gamma_{R_2,S_2}(b_\star)\neq0\), \(w_2(0)\) admits no lift in
\(\mathcal{J}_\tau^{(2)}\), so \(w_2(0)\notin J_\tau^2\), which forces
\(J_\tau^2=0\) (matching \cite[Theorem~C(e)]{Astashkevich1997}).

\emph{Step 4: Taking \(\tau\to0\).}
Diagonalize the symmetric matrix \(\bar B_{V,\tau}\) on the free
\(\mathbb{C}[[\tau]]\)-module \(\mathcal{V}_N\) by congruence
\(U(\tau)^{\mathsf{T}}\bar B_{V,\tau}U(\tau)
=\operatorname{diag}(\tau^{e_1}u_1(\tau),\ldots,\tau^{e_m}u_m(\tau))\) with
\(u_a(0)\neq0\).  Since \(J_\tau^2=0\), every exponent \(e_a\) lies in
\(\{0,1\}\), with \(\{v_a(0):e_a=0\}\) descending to a basis of
\(L(c_\star,\mu)_N\) and \(\{v_a(0):e_a=1\}\) forming a basis of
\(L(c_\star,\nu_2)_{N-R_2S_2}\).  For \(a(\tau),b(\tau)\in\mathcal{J}_\tau^{(1)}\)
homogeneous of levels \(N_a,N_b\), we have
\(\bar\rho_{V,\tau}(a(\tau),b(\tau))=O(\tau)\) (because
\(J_\tau^1=L(c_\star,\nu_2)\subset\rad\bar\rho_{V,0}\)).  Shifting either lift
by \(\tau u(\tau)\) changes \(\bar\rho_{V,\tau}(a(\tau),b(\tau))\) by
\(O(\tau^2)\) (again because \(J_\tau^1\subset\rad\bar\rho_{V,0}\)), so
\(\bar\rho^{(1)}_\tau(a(0),b(0)):=[\tau^1]\bar\rho_{V,\tau}(a(\tau),b(\tau))\)
is well-defined on \(L(c_\star,\nu_2)\times L(c_\star,\nu_2)\).  Taking
\([\tau^1]\) of the Ward identity
\(\bar\rho_{V,\tau}(L_{-n}a(\tau),b(\tau))-\bar\rho_{V,\tau}(a(\tau),L_nb(\tau))
=(N_a-N_b+nd(\tau))\bar\rho_{V,\tau}(a(\tau),b(\tau))\)
drops the \(nd'(0)\tau\cdot O(\tau)=O(\tau^2)\) term and shows that
\(\bar\rho^{(1)}_\tau\) satisfies the exact \(V_{d_\star}(1)\) Ward identity
on \(L(c_\star,\nu_2)\).  By uniqueness of the contravariant and primary forms
on the irreducible quotient \(L(c_\star,\nu_2)\), we have
\(\bar B^{(1)}_\tau=\alpha\,\gamma_{R_2,S_2}(b_\star)\,\bar B_{L(\nu_2)}\) and
\(\bar\rho^{(1)}_\tau=\alpha\,\tilde g_\rho(t_0,\mu)\,\bar\rho_{L(\nu_2),d_\star}\),
whose ratio is
\[
 \frac{\alpha\,\tilde g_\rho(t_0,\mu)}{\alpha\,\gamma_{R_2,S_2}(b_\star)}
 =\frac{\tilde g_\rho(t_0,\mu)}{\gamma_{R_2,S_2}(b_\star)}
 =c_{R_2,S_2}(t_0;u,v)
 =\Omega_{u,v}(\mu\to\nu_2)
\]
(specializing the rational identity of Theorem~\ref{thm:first-layer-ratio} at
\(t=t_0\), where both \(\gamma_{R_2,S_2}(b_\star)\) and the reduced box
denominator of Proposition~\ref{prop:bgg-cocycle}(i) are nonzero).
Therefore the \(\tau^1\) factors cancel in the diagonalized trace on
\(\mathcal{V}_N\), giving
\[
 \lim_{\tau\to0}\Tr_{V_{R_1,S_1}(b(\tau))_N}\bigl(V_{d(\tau)}(1)\bigr)
 =T^{\mathrm{irr}}_N(c_\star,\mu,d_\star)
  +\Omega_{u,v}(\mu\to\nu_2)\,T^{\mathrm{irr}}_{N-R_2S_2}(c_\star,\nu_2,d_\star).
\]
Meanwhile, by Lemma~\ref{lem:joint-holomorphicity} and
Proposition~\ref{prop:bgg-cocycle}(i), the remaining terms in
\eqref{eq:generic-tau-split} are regular at \(\tau=0\) and converge to
\(F_N(c_\star,\mu,d_\star)\) and
\(\Omega_{u,v}(\mu\to\nu_1)F_{N-R_1S_1}(c_\star,\nu_1,d_\star)\).  Taking
\(\tau\to0\) in \eqref{eq:generic-tau-split} yields \eqref{eq:two-term-id}.
\end{proof}

\begin{proof}[Proof of Theorem~\ref{thm:gen-rocha-caridi}]
For each \(\mu\in\mathcal{O}_{p,p'}(r,s)\), define the rescaled formal power
series
\[
 \widetilde{\mathcal{T}}(\mu)
 :=\Omega_{u,v}(\lambda\to\mu)\,q^{\mu-\lambda}\,
   \mathcal{T}^{\mathrm{irr}}_{c_\star,\mu,d}(q),
 \qquad
 \widetilde{\mathcal{F}}(\mu)
 :=\Omega_{u,v}(\lambda\to\mu)\,q^{\mu-\lambda}\,
   \mathcal{F}^{\mathrm{Verma}}_{c_\star,\mu,d}(q),
\]
with \(\Omega_{u,v}(\lambda\to\lambda):=1\).  By the cocycle property
\(\Omega_{u,v}(\lambda\to\mu)\Omega_{u,v}(\mu\to\nu)=\Omega_{u,v}(\lambda\to\nu)\)
(Proposition~\ref{prop:bgg-cocycle}(ii)), multiplying \eqref{eq:two-term-id}
by \(\Omega_{u,v}(\lambda\to\mu)q^{\mu-\lambda}\) yields
\begin{equation}\label{eq:scaled-two-term}
 \widetilde{\mathcal{T}}(\mu)+\widetilde{\mathcal{T}}(\nu_2)
 =\widetilde{\mathcal{F}}(\mu)-\widetilde{\mathcal{F}}(\nu_1)
\end{equation}
for every \(\mu\in\mathcal{O}_{p,p'}(r,s)\) with primitive children
\(\{\nu_1,\nu_2\}\).  Now apply \eqref{eq:scaled-two-term} along the alternating
chain \(\mu_0:=\lambda\), \(\mu_1:=\mu^-_1\), \(\mu_2:=\mu^+_2\),
\(\mu_3:=\mu^-_3\), \(\ldots\), choosing at step \(j\geq0\) the child
\(\nu_1=\mu^{\epsilon_j}_{j+1}\) and \(\nu_2=\mu^{-\epsilon_j}_{j+1}\) with
\(\epsilon_j=(-1)^j\) (so \(\nu_2=\mu_{j+1}\) for every \(j\geq0\)):
\[
 \widetilde{\mathcal{T}}(\mu_j)+\widetilde{\mathcal{T}}(\mu_{j+1})
 =\widetilde{\mathcal{F}}(\mu_j)-\widetilde{\mathcal{F}}(\mu^{(-1)^j}_{j+1})
 \qquad (j\geq0).
\]
For any fixed level \(N\geq0\), choose \(K\) so large that
\(\min(\ell^+_{K+1},\ell^-_{K+1})>N\).  Taking the finite alternating sum
\(\sum_{j=0}^K(-1)^j\) telescopes the left-hand side to
\(\widetilde{\mathcal{T}}(\lambda)+(-1)^K\widetilde{\mathcal{T}}(\mu_{K+1})
\equiv\mathcal{T}^{\mathrm{irr}}_{c_\star,\lambda,d}(q)\pmod{q^{N+1}}\),
while the right-hand side equals
\[
 \widetilde{\mathcal{F}}(\lambda)
 +\sum_{j=1}^K(-1)^j\bigl(\widetilde{\mathcal{F}}(\mu^+_j)
                         +\widetilde{\mathcal{F}}(\mu^-_j)\bigr)
 +(-1)^{K+1}\widetilde{\mathcal{F}}(\mu^{(-1)^K}_{K+1}).
\]
Since \(\ell^\pm_{K+1}>N\), the final term vanishes modulo \(q^{N+1}\).  As
\(N\geq0\) is arbitrary, this proves \eqref{eq:gen-rocha-caridi}.

Part~(a) is immediate: when \(k=m=0\) we have \(d=0\),
\(\Omega_{1,1}\equiv1\), and \(F_N(c_\star,\mu,0)=p(N)\) for all \(\mu\).
For part~(b), \cite[Theorem~4.3]{CompanionTorus} computes the exact finite
partial-fraction expansion of \(H_M(c_\star,\Delta,d)\) in \(\Delta\) from the
generic Poghossian--HJS recursion with no numerical \(\varepsilon\)-limit; by
Theorem~\ref{thm:descent-regularity}, all poles at
\(\Delta=\mu\in\mathcal{O}_{p,p'}(r,s)\) have identically zero Laurent
coefficients, so omitting the identically zero principal part at \(\Delta=\mu\)
lets \(H_M(c_\star,\mu,d)\) be evaluated by direct substitution of
\(\Delta=\mu\) into the non-\(\mu\) partial-fraction summands, converting to
\(F_M(c_\star,\mu,d)=\sum_{i=0}^Mp(i)H_{M-i}(c_\star,\mu,d)\) and summing the
finite truncation of \eqref{eq:gen-rocha-caridi} with \(\ell^\pm_j\leq N\).

For part~(c), apply \eqref{eq:gen-rocha-caridi} with \(\lambda\) replaced by
any \(\mu\in\mathcal{O}_{p,p'}(r,s)\) (which holds by the same telescoping
argument starting from \(\mu\)): in matrix form on the ordered orbit
\(\mathcal{O}_{p,p'}(r,s)\), this expresses the vector
\((\widetilde{\mathcal{T}}(\mu))_\mu\) as \(U_0^{-1}(\widetilde{\mathcal{F}}(\mu))_\mu\),
where \(U_0\) is the unitriangular zeta matrix with \((U_0)_{\mu,\nu}=1\) for
\(\nu\succeq\mu\) in the braided poset.  Indeed, for any interval
\([\mu,\nu]\) of layer distance \(d>0\), the braided poset contains one bottom
vertex \(\mu\), two vertices at each intermediate distance \(1\leq i\leq d-1\),
and one top vertex \(\nu\), giving
\(1+2\sum_{i=1}^{d-1}(-1)^i+(-1)^d=0\); thus
\((U_0^{-1})_{\mu,\nu}=(-1)^{\mathrm{dist}(\mu,\nu)}\).
Multiplying by \(U_0\) recovers
\(\widetilde{\mathcal{F}}(\lambda)=\sum_{\nu\succeq\lambda}\widetilde{\mathcal{T}}(\nu)\),
i.e.
\[
 F_N(c_\star,\lambda,d)
 =T^{\mathrm{irr}}_N(c_\star,\lambda,d)
  +\sum_{j=1}^{j_N(\lambda)}\sum_{\pm}
    \Omega_{u,v}(\lambda\to\mu^\pm_j)\,
    T^{\mathrm{irr}}_{N-\ell^\pm_j}(c_\star,\mu^\pm_j,d).
\]
Every orbit level \(\ell^\pm_j\) (\(j\geq1\)) is pairwise distinct from all
others: each layer-\((j+1)\) weight strictly exceeds both layer-\(j\) weights
by positive primitive arrow levels, and within each layer \(\ell^+_j\neq\ell^-_j\)
by Lemma~\ref{lem:bgg-arrows}(iv).  Comparing this expansion with
\eqref{eq:jantzen-trace-decomp} inductively in increasing order of \(\ell^\pm_j\)
(using \(T^{\mathrm{irr}}_0=1\)) proves
\(\sigma^\pm_j(\lambda)/\gamma^\pm_j(\lambda)=\Omega_{u,v}(\lambda\to\mu^\pm_j)\)
for every \(j\geq1\).
\end{proof}

\begin{proof}[Proof of Corollary~\ref{cor:collision-free}]
By \(p\leftrightarrow p'\) duality it suffices to treat
\(d(b)=h_{1,2m+1}(b)\) with \(0\leq m<\min(s,p'-s)\).  As noted in the proof of
Lemma~\ref{lem:joint-holomorphicity}, \(R_{R,S}(b,h_{1,2m+1}(b))\equiv0\)
unless \(S\leq m\) (since \(k=0\) precludes \(R\leq k\)).  Notice that
\(2m\leq 2\min(s,p'-s)-2\leq p'-2<p'\) (for both even and odd \(p'\)).  Now
check the three obstructions for active labels \((R,S)\) with \(1\leq S\leq m\):
\begin{enumerate}
 \item \emph{No \(b\)-poles in \(R_{R,S}(b,d(b))\) at \(b=b_\star\):}
   the denominator factors of \(A_{R,S}(b)\) are \(ab+b'b^{-1}=(at+b')/b\) with
   \(1-S\leq b'\leq S\) and \((a,b')\neq(0,0)\).  At \(t_\star=-p'/p\),
   \(pb'-p'a=0\) requires \(p'\mid b'\).  Since \(|b'|\leq S\leq m<p'/2\), this
   forces \(b'=0\) and hence \(a=0\), which is excluded.
 \item \emph{No primary collisions:}
   if \(h_{R_1,S_1}(b_\star)=h_{R_2,S_2}(b_\star)\) with
   \(1\leq S_1,S_2\leq m\), then \(p'R_1-pS_1=\pm(p'R_2-pS_2)\) implies
   \(p'\mid(S_1\mp S_2)\).  Since \(|S_1-S_2|<p'/2\) and
   \(0<S_1+S_2\leq 2m<p'\), we must have \(S_1=S_2\) and the \(+\) sign, hence
   \(R_1=R_2\).
 \item \emph{No secondary collisions:}
   if \(h_{R_1,-S_1}(b_\star)=h_{R_2,S_2}(b_\star)\) with
   \(1\leq S_1,S_2\leq m\), then \(p'R_1+pS_1=\pm(p'R_2-pS_2)\).  The \(-\)
   sign gives \(p'(R_1+R_2)=p(S_2-S_1)\), forcing \(p'\mid(S_2-S_1)\), hence
   \(S_1=S_2\) and \(R_1+R_2=0\), impossible.  The \(+\) sign gives
   \(p'(R_2-R_1)=p(S_1+S_2)\), forcing \(p'\mid(S_1+S_2)\), impossible since
   \(0<S_1+S_2\leq 2m<p'\).
\end{enumerate}
Finally, by Lemma~\ref{lem:joint-holomorphicity}, no active pole
\(h_{R,S}(b_\star)\) coincides with any BGG weight
\(\mu\in\mathcal{O}_{p,p'}(r,s)\).  Therefore every denominator in the
truncated recursion for \(H_N(c_\star,\mu,d_\star)\) is nonzero at
\(b=b_\star\).
\end{proof}

\section{Minimal-Model Examples and Exact Verification}\label{sec:examples}

We illustrate Theorems~\ref{thm:descent-regularity}--\ref{thm:gen-rocha-caridi}
and Corollary~\ref{cor:collision-free} on three minimal models.  All data below
are reproduced in exact rational and symbolic arithmetic by the script
\texttt{audit\_irreducible\_torus.py} in \texttt{irreducible/tools/}.

\subsection{Ising model \texorpdfstring{\((3,4)\), \(\lambda=1/16\), \(d=1/2\)}{(3,4), lambda=1/16, d=1/2}}

Here \(c_\star=1/2\), \(t_\star=-4/3\), \((r,s)=(1,2)\), and
\(\mathcal{F}_{3,4}(1,2)=\{h_{1,1},h_{1,3}\}=\{0,1/2\}\).  For \(d=h_{1,3}=1/2\)
(\(k=0,m=1\)), the first seven BGG weights and rational cocycle coefficients
\((-1)^j\Omega_{1,3}(\lambda\to\mu^\pm_j)\) through level \(30\) (crossing
Jantzen layers \(0,1,2,3\)) are
\[
 \begin{array}{c|ccccccc}
  \ell & 0 & 2 & 4 & 10 & 14 & 24 & 30 \\\hline
  (-1)^j\Omega_{1,3} & 1 & -\dfrac{15}{7} & \dfrac{195}{77} & -\dfrac{663}{209} & \dfrac{16575}{4807} & -\dfrac{160225}{40641} & \dfrac{1185665}{284487}
 \end{array}
\]
By Corollary~\ref{cor:collision-free}, the truncated recursion for
\(H_N(1/2,\mu,1/2)\) retains only labels \((R,1)\) (\(R\geq1\)) and is
completely pole-free and collision-free at \(b=b_\star\).  Evaluating
\eqref{eq:gen-rocha-caridi} cancels every denominator prime \(7,11,17,19,23,\ldots\)
and agrees identically for all checked levels \(0\leq N\leq 30\) with the Euler
pentagonal series \(\phi(q)=\prod_{n=1}^\infty(1-q^n)\) (matching
\(q^{1/16-1/48}\phi(q)=\eta(\tau)\) through \(N\leq30\), and agreeing with
direct quotient-Gram inversion for \(0\leq N\leq 8\)).

\subsection{Lee--Yang model \texorpdfstring{\((2,5)\), \(\lambda=-1/5\), \(d=-1/5\)}{(2,5), lambda=-1/5, d=-1/5}}

Here \(c_\star=-22/5\), \(t_\star=-5/2\), \((r,s)=(1,2)\), and
\(\mathcal{F}_{2,5}(1,2)=\{0,-1/5\}\) with \(d=-1/5=h_{1,3}(b_\star)\)
(\(k=0,m=1\)).  The BGG levels and weights through level \(25\) (layers
\(0,1,2,3\)) are
\[
 \begin{array}{c|ccccccc}
  \ell & 0 & 2 & 3 & 9 & 11 & 21 & 24 \\\hline
  (-1)^j\Omega_{1,3} & 1 & -\dfrac{3}{7} & -\dfrac{2}{3} & \dfrac{44}{119} & \dfrac{78}{133} & -\dfrac{52}{153} & -\dfrac{299}{551}
 \end{array}
\]
Direct quotient-Gram inversion matches \eqref{eq:gen-rocha-caridi} in exact
rational arithmetic for all \(0\leq N\leq 9\) (explicitly crossing Jantzen
layer \(2\) at \(N=9\), where \(F_9\) has a \(10\)-digit denominator containing
\(7,11,13,17,\ldots\) and the BGG sum collapses to
\(T^{\mathrm{irr}}_9=27946402/5^{10}\)).  Evaluating
\eqref{eq:gen-rocha-caridi} via Corollary~\ref{cor:collision-free} for
\(0\leq N\leq 25\) shows by exact calculation that every coefficient
\(T^{\mathrm{irr}}_N\) in this checked range has a pure power of \(5\) as its
denominator, starting with
\[
 1,\;\frac{2}{5},\;\frac{17}{25},\;\frac{98}{125},\;\frac{714}{625},\;
 \frac{18768}{15625},\;\frac{140231}{78125},\;\frac{721786}{390625},\;
 \frac{4990044}{1953125},\;\ldots
\]
By contrast, the reflected weight \(d=h_{1,-3}(b_\star)=14/5\in\mathcal{F}^\pm_{2,5}(1,2)\setminus\mathcal{F}_{2,5}(1,2)\)
illustrates the strict inclusion in Theorem~\ref{thm:descent-regularity}:
\(F_N(-22/5,\Delta,14/5)\) is regular at \(\Delta=-1/5\) for every \(N\geq0\),
yet \(\rho_2(-22/5,-1/5,14/5)\) fails to descend to \(L_2(-22/5,-1/5)\).

\subsection{Tricritical Ising model \texorpdfstring{\((4,5)\), \(\lambda=3/80\)}{(4,5), lambda=3/80}, two-index fusion}

For \((p,p')=(4,5)\) (\(c_\star=7/10\), \(t_\star=-5/4\)) at
\((r,s)=(2,2)\) (\(\lambda=h_{2,2}=3/80\), \(\ell^+_1=4\), \(\ell^-_1=6\)),
the fusion set \eqref{eq:fusion-set} has four elements:
\(\mathcal{F}_{4,5}(2,2)=\{h_{1,1},h_{1,3},h_{3,1},h_{3,3}\}=\{0,\frac{3}{5},\frac{3}{2},\frac{1}{10}\}\).
For \(d=h_{3,3}=1/10\) (\(k=m=1\)), active labels \((1,4)\) and \((3,1)\)
collide off-orbit at \(b=b_\star\), exercising the full two-index case of
Theorem~\ref{thm:gen-rocha-caridi}.  The first-layer BGG weights
\((-1)\Omega_{u,v}(\lambda\to\mu^\pm_1)\) at levels \((4,6)\) are
\((17/3,-154/39)\) for \(d=3/5\), \((-1463/39,33649/663)\) for \(d=3/2\), and
\((77/39,32186/25857)\) for \(d=1/10\); all three match direct quotient-Gram
inversion identically in \(\mathbb{Q}\) for \(0\leq N\leq 7\).

\subsection{Audit summary}

The standalone script \texttt{irreducible/tools/audit\_irreducible\_torus.py}
independently checks: (1)~Theorem~\ref{thm:descent-regularity} across allowed,
forbidden, and both types of half-fusion weights (including \((3,5)\),
\((r,s)=(1,3)\), \(d=h_{1,5}=8/5\), where \(P_{1,3}(d)=0\) gives
\(\ord_xF_3=0\) while \(P_{2,2}(d)\neq0\) produces \(\ord_xF_4=-1\));
(2)~Theorem~\ref{thm:first-layer-ratio} directly from the two-weight Ward
identity and symbolic Gamma simplification in \(\mathbb{Q}(t)\);
(3)~Theorem~\ref{thm:gen-rocha-caridi} against exact \(\mathbb{Q}\)
quotient-Gram traces across Lee--Yang (\(N\leq9\)), Ising (\(N\leq8\)), and
all three nonzero Tricritical Ising fusion insertions (\(N\leq7\)); and
(4)~Corollary~\ref{cor:collision-free} for Ising (\(N\leq30\)) and Lee--Yang
(\(N\leq25\)).

\section*{AI-use disclosure}

AI systems had a substantial role in developing and checking the proofs,
writing the exact verification scripts, and drafting the manuscript.  The
author directed the research and takes responsibility for all mathematical
claims, citations, computations, and wording.

\section*{Code availability}

The exact rational and symbolic verification script in
\texttt{irreducible/tools/} reproduces every verification reported in
Section~\ref{sec:examples}.  The repository is at
\href{https://github.com/skypher/virasoro}{github.com/skypher/virasoro}.

\begin{thebibliography}{99}

\bibitem{Astashkevich1997}
A.~Astashkevich,
\emph{On the structure of Verma modules over Virasoro and Neveu--Schwarz algebras},
Comm.\ Math.\ Phys.\ \textbf{186} (1997), 531--562,
arXiv:hep-th/9511032.

\bibitem{Felder1989}
G.~Felder,
\emph{BRST approach to minimal models},
Nucl.\ Phys.\ B \textbf{317} (1989), 215--236.

\bibitem{FelderSilvotti1989}
G.~Felder and R.~Silvotti,
\emph{Free field representation of minimal models on a Riemann surface},
Phys.\ Lett.\ B \textbf{231} (1989), 411--416.

\bibitem{HJS2010}
L.~Hadasz, Z.~Jask\'olski, and P.~Suchanek,
\emph{Recursive representation of the torus 1-point conformal block},
J.\ High Energy Phys.\ \textbf{2010}, no.~1, 063,
arXiv:0911.2353.

\bibitem{IoharaKoga2011}
K.~Iohara and Y.~Koga,
\emph{Representation Theory of the Virasoro Algebra},
Springer Monographs in Mathematics, Springer, London, 2011.

\bibitem{JSF2018}
N.~Javerzat, R.~Santachiara, and O.~Foda,
\emph{Notes on the solutions of Zamolodchikov-type recursion relations in Virasoro minimal models},
J.\ High Energy Phys.\ \textbf{2018}, no.~8, 183,
arXiv:1806.02790.

\bibitem{Poghossian2009}
R.~Poghossian,
\emph{Recursion relations in CFT and \(N=2\) SYM theory},
J.\ High Energy Phys.\ \textbf{2009}, no.~12, 038,
arXiv:0909.3412.

\bibitem{Ribault2019}
S.~Ribault,
\emph{On 2d CFTs that interpolate between minimal models},
SciPost Phys.\ \textbf{6} (2019), 075,
arXiv:1809.03722.

\bibitem{Ribault2026}
S.~Ribault,
\emph{Exactly solvable conformal field theories},
SciPost Phys.\ Lect.\ Notes \textbf{130} (2026),
arXiv:2411.17262v4.

\bibitem{RochaCaridi1985}
A.~Rocha-Caridi,
\emph{Vacuum vector representations of the Virasoro algebra},
in \emph{Vertex Operators in Mathematics and Physics},
MSRI Publ.\ \textbf{3}, Springer, New York, 1985, pp.~451--473.

\bibitem{CompanionTorus}
L.~P.~Polzer,
\emph{Exact algebraic confluence for torus one-point Virasoro-block recursions},
companion manuscript in \href{https://github.com/skypher/virasoro}{github.com/skypher/virasoro} (2026).

\bibitem{Zamolodchikov2004}
Al.~B.~Zamolodchikov,
\emph{Higher equations of motion in Liouville field theory},
Int.\ J.\ Mod.\ Phys.\ A \textbf{19S2} (2004), 510--523,
arXiv:hep-th/0312279.

\end{thebibliography}

\end{document}
